\documentclass[11pt]{amsart}
\usepackage[T1]{fontenc}
\usepackage[utf8]{inputenc}
\usepackage{lmodern,amssymb,mathtools,enumitem,booktabs,array}
\usepackage[margin=1.15in]{geometry}
\usepackage[colorlinks=true,linkcolor=blue,citecolor=blue,urlcolor=blue]{hyperref}
\newtheorem{theorem}{Theorem}[section]
\newtheorem{lemma}[theorem]{Lemma}
\newtheorem{proposition}[theorem]{Proposition}
\newtheorem{corollary}[theorem]{Corollary}
\theoremstyle{definition}
\newtheorem{definition}[theorem]{Definition}
\newtheorem{construction}[theorem]{Construction}
\newtheorem{example}[theorem]{Example}
\theoremstyle{remark}
\newtheorem{remark}[theorem]{Remark}
\newcommand{\N}{\mathbb{N}}
\newcommand{\R}{\mathbb{R}}
\newcommand{\cD}{\mathcal{D}}
\newcommand{\cC}{\mathcal{C}}
\newcommand{\cS}{\mathcal{S}}
\newcommand{\cE}{\mathcal{E}}
\newcommand{\cK}{\mathcal{K}}
\newcommand{\cT}{\mathcal{T}}
\newcommand{\fM}{\mathfrak{M}}
\newcommand{\diam}{\operatorname{diam}}
\newcommand{\cl}{\operatorname{cl}}
\newcommand{\dom}{\operatorname{dom}}
\newcommand{\cone}{\mathsf{c}}
\newcommand{\SHE}{\mathit{SHE}}
\newcommand{\HE}{\mathit{HE}}
\numberwithin{equation}{section}

\begin{document}
\title[Three questions on mosaic spaces]{Mosaic spaces generated by continua and by Peano continua: answers to three questions recorded by Lavallee}
% Author name and e-mail as supplied by the author. No affiliation, ORCID,
% funding statement or author-responsibility statement has been added.
\author{Lucas Alves de Almeida}
\email{Lucasderiva8x@gmail.com}
\subjclass[2020]{Primary 54D55; Secondary 54D05, 54F15, 54G20}
\keywords{mosaic space, sequential space, $m$-continuum space, curve space, path-determined space, local arcwise connectedness, continuum-wise connectedness, cone}

\begin{abstract}
Lavallee (1964) recorded three questions about Davison's mosaic spaces: whether a locally ($m$-strongly) connected $m$-continuum space must be a quasi-curve space; whether the converse of her theorem on open subspaces of $m$-continuum spaces holds; and Davison's question whether a locally arcwise connected mosaic space must be a quasi-curve space. We show that all three have negative answers. Davison's question is answered by published examples of Franklin and Smith-Thomas (1979) and of van Douwen (1981), once ``path-determined'' is identified with ``curve space'', and an example sketched by Curtis (1980) bears on the first question. We give self-contained constructions and proofs: a metric space without arcs, a completely regular space in which no metric continuum captures a certain convergent sequence, and a cone construction that answers the second question. We also prove affirmative results for first countable spaces. We claim no priority for these answers.
\end{abstract}
\maketitle

\section{Introduction}\label{sec:intro}

\subsection{Mosaic spaces}
Many familiar spaces are assembled from compact metrizable pieces in such a way that a set is closed exactly when it meets every piece in a closed set. A metric space is assembled in this way from its convergent sequences, a CW complex from its finite subcomplexes. Davison \cite{davison59} made this the definition of a class of spaces: a \emph{mosaic} is a family of compact metric spaces whose underlying sets cover a set $X$ and whose topologies are compatible on overlaps, and the \emph{mosaic space} is $X$ with the largest topology inducing the given one on every piece. Precise definitions are in Section~\ref{sec:prelim}. According to Lavallee \cite[Theorem~1.6]{lavallee64}, Davison showed that mosaic spaces are exactly the spaces whose topology is ``full'' and in which sequences have unique limits; in the terminology introduced shortly afterwards by Franklin \cite{franklin65}, they are the sequential spaces with unique sequential limits. In particular every Hausdorff sequential space is a mosaic space.

The pieces of a mosaic can be restricted. Lavallee \cite{lavallee64} studied three cases: the pieces are \emph{metric continua} (compact, connected, metrizable), giving the \emph{$m$-continuum spaces}; the pieces are compact, locally connected and metrizable, giving the \emph{quasi-curve spaces}; the pieces are Peano continua, giving Davison's \emph{curve spaces}. She proved that the last two classes coincide \cite[Theorem~6.1]{lavallee64}. In terms of sequences (Corollary~\ref{cor:classes}), a mosaic space is an $m$-continuum space when every convergent sequence has a subsequence which lies, together with its limit, in one metric continuum inside the space, and it is a quasi-curve space when that continuum can always be chosen locally connected, that is, when the subsequence and its limit lie on the image of a path.

These are \emph{global} conditions on convergent sequences: one compact metrizable connected set has to contain the limit $x$ and infinitely many terms $x_n$. The natural \emph{local} conditions say something different. Local $m$-strong connectedness (every point has arbitrarily small open neighbourhoods in which any two points lie in a common metric continuum) provides, inside any given neighbourhood of $x$, metric continua joining $x$ to each of the terms $x_n$ that lie in a suitable smaller neighbourhood; local arcwise connectedness provides arcs in the same way. Local connectedness provides less: connected open neighbourhoods, but not necessarily any metric continuum joining $x$ to a nearby point.

To pass from such connecting continua $K_n\ni x,x_n$ to a single continuum containing $x$ and infinitely many $x_n$, one has to choose them so that their union is compact and metrizable. A sufficient way to obtain a compact union is to choose the connecting continua so that every neighbourhood of $x$ contains all but finitely many of them (Corollary~\ref{cor:image}(3)). If $x$ has a countable neighbourhood base, this can be arranged by a diagonal choice (Section~\ref{sec:fc}). Without one, the continua supplied by the local condition depend on the neighbourhood, and nothing guarantees that a single sequence of them has compact metrizable union; the examples below show that capture can indeed fail. There is a second, independent gap: connecting sets that are continua need not contain arcs. In a completely metrizable space local connectedness does produce arcs, by a classical theorem of Mazurkiewicz, Moore and Menger \cite[Section~6.3]{engelking}, but in an incomplete metric space it need not.

\subsection{The questions}
The following questions appear together on the last page of Lavallee's article \cite[p.~1333]{lavallee64}, after her Theorem~5.2, which states that if $X$ and all its open subspaces are $m$-continuum spaces then $X$ is locally $m$-strongly connected.

\begin{enumerate}[label=\textbf{Q\arabic*.},leftmargin=3.2em]
\item Lavallee writes that it would be of interest to know whether ``a locally ($m$-strongly) connected $m$-continuum space is a quasi-curve space''. The parenthesis allows two readings, and we treat both:
\begin{itemize}[leftmargin=2.2em]
\item[\textbf{Q1a.}] Is every locally connected $m$-continuum space a quasi-curve space?
\item[\textbf{Q1b.}] Is every locally $m$-strongly connected $m$-continuum space a quasi-curve space?
\end{itemize}
\item In the same passage she asks whether the converse of Theorem~5.2 holds: if $X$ is a locally $m$-strongly connected $m$-continuum space, is every open subspace of $X$ an $m$-continuum space?
\item She reports a question put to her by Davison, and describes it as unanswered: is every locally arcwise connected mosaic space a quasi-curve space?
\end{enumerate}

A metric continuum is connected, so local $m$-strong connectedness implies local connectedness. Thus the hypothesis of Q1b is formally stronger than that of Q1a, an affirmative answer to Q1a would imply one to Q1b, and a counterexample to Q1b is a counterexample to Q1a. The two hypotheses are genuinely different, even for $m$-continuum spaces (Example~\ref{ex:cZ}).

Lavallee's article gives no further motivation for these questions beyond their position after Theorem~5.2 and Corollary~6.2 (quasi-curve spaces are $m$-continuum spaces), and we do not attribute one to her or to Davison. Our own reading is that they ask whether the hierarchy
\[\text{curve space}=\text{quasi-curve space}\ \Longrightarrow\ \text{$m$-continuum space}\ \Longrightarrow\ \text{mosaic space}\]
collapses in the presence of the corresponding local connectedness property. The answer in each case is that it does not, and the examples identify the two gaps described above as the reasons.

\subsection{Answers}
Table~\ref{tab:answers} summarises the situation; the terms used are defined in Section~\ref{sec:prelim}.

\begin{table}[ht]
\renewcommand{\arraystretch}{1.25}
\small\setlength{\tabcolsep}{4pt}\begin{tabular}{@{}p{0.09\textwidth}p{0.075\textwidth}p{0.345\textwidth}p{0.365\textwidth}@{}}
\toprule
Question & Answer & Published antecedent & In this article\\
\midrule
Q1a, Q1b & No & Curtis \cite[Example~1.8]{curtis80}: a plane set, described with a figure, which is locally continuum-connected and not locally path connected. Not related there to Lavallee's question. & Theorem~\ref{thm:arcless} and Corollary~\ref{cor:Q1}: a metric space with no arcs in which all balls are continuum-wise connected. Example~\ref{ex:cZ} for Q1a alone.\\
Q2 & No & None found in the sources examined. & Theorem~\ref{thm:conetransfer} and Corollary~\ref{cor:Q2}: the cone over the space $G$. Affirmative for first countable spaces: Proposition~\ref{prop:fc}.\\
Q3 & No & Franklin and Smith-Thomas \cite[Section~3]{fst79}; van Douwen \cite{vandouwen81}. Neither mentions mosaics. & Proposition~\ref{prop:pathdet} (translation of terminology); Theorem~\ref{thm:G} and Corollary~\ref{cor:Q3}: the space $G$. Affirmative for first countable spaces: Proposition~\ref{prop:fcpath}.\\
\bottomrule
\end{tabular}
\caption{The questions and their answers.}\label{tab:answers}
\end{table}

Our principal results are the following.

\begin{itemize}[leftmargin=2em]
\item \emph{A metric space without arcs.} For every metric space $Y$ with at least two points that contains no arc there is a connected metric space $E(Y)\supseteq Y$ which contains no arc and in which any two points $x,y$ lie in a continuum of diameter $d(x,y)$ (Theorem~\ref{thm:arcless}). Every open subspace of $E(Y)$ is an $m$-continuum space, and $E(Y)$ is not a quasi-curve space. This answers Q1a and Q1b.
\item \emph{The space $G$.} There is a completely regular, locally arcwise connected mosaic space $G$ with a sequence $x_n\to\infty$ such that no closed connected subset of $G$ with countable spread, in particular no metric continuum, contains $\infty$ and infinitely many $x_n$ (Theorem~\ref{thm:G}). This answers Q3, and $G$ is not even an $m$-continuum space.
\item \emph{The published examples.} For Hausdorff spaces, van Douwen's ``path-determined'' coincides with ``curve space'' (Proposition~\ref{prop:pathdet}), so the spaces of \cite{fst79} and \cite{vandouwen81} answer Q3. We prove in addition that no compact connected subset of van Douwen's space contains infinitely many of its distinguished integer points (Proposition~\ref{prop:SHE}), so that space is not an $m$-continuum space either.
\item \emph{Cones.} The cone $\cone Y$ over a mosaic space $Y$ is a contractible $m$-continuum space; it is locally path connected if $Y$ is; and every open subspace of $\cone Y$ is an $m$-continuum space if and only if the same holds for $Y$ (Theorem~\ref{thm:conetransfer}). Consequently $\cone G$ answers Q2.
\item \emph{First countable spaces and the comparison of the conditions.} Q2 and Q3 have affirmative answers for first countable mosaic spaces, in a form that captures a whole tail of the sequence (Propositions~\ref{prop:fc} and \ref{prop:fcpath}). Among $m$-continuum spaces the conditions ``quasi-curve'', ``every open subspace is an $m$-continuum space'', ``locally $m$-strongly connected'' and ``locally connected'' form a chain of implications none of which can be reversed, and local arcwise connectedness is independent of the second condition (Theorem~\ref{thm:chain}).
\end{itemize}

\subsection{Attribution and novelty}\label{sec:attribution}
We make no claim of priority for the negative answers or for the phenomena exhibited by these examples. Our purpose is to connect the historical questions with relevant published constructions and to provide explicit proofs and further deductions. For certain statements proved below we did not locate an earlier formulation in the sources we examined; this observation does not establish novelty, and the search was not exhaustive. In detail:

\begin{enumerate}[label=(\arabic*),leftmargin=2em]
\item \emph{Proved in the literature.} The basic theory of mosaic spaces is due to Davison and Lavallee; the facts we use are restated with proofs in Section~\ref{sec:prelim} to make the article self-contained (see Section~\ref{sec:sources}). Franklin and Smith-Thomas \cite{fst79} constructed a Hausdorff, sequential, locally path connected space that is not path-determined, and van Douwen \cite{vandouwen81} a regular, Fr\'echet one. That first countable, locally path connected spaces are path-determined is stated in \cite{vandouwen81} as taken from \cite{fst79}. Curtis \cite{curtis80} exhibited local continuum-connectedness without local path connectedness in a subset of the plane.
\item \emph{Published results whose relevance to Lavallee's questions is established here.} That the spaces of \cite{fst79} and \cite{vandouwen81} answer Davison's question (Section~\ref{sec:published}), and that Curtis's example would answer Q1 (Remark~\ref{rem:curtis}).
\item \emph{Alternative constructions and self-contained verifications.} The space $E(Y)$ of Section~\ref{sec:arcless} exhibits the phenomenon of Curtis's example in a stronger form (no arcs) with complete proofs. The space $G$ of Section~\ref{sec:G} was designed from the two ingredients of \cite{fst79}, Arens' space and a maximal almost disjoint family; its mechanism is essentially that of \cite{vandouwen81}, and we claim no novelty for it. What distinguishes $G$ is its presentation as a mosaic of explicit compact metric pieces, which makes every verification short.
\item \emph{Further deductions.} The obstruction to capturing a sequence by \emph{continua}, as opposed to paths, for $G$ (Proposition~\ref{prop:nocapture}) and for van Douwen's space (Proposition~\ref{prop:SHE}); the complete regularity of $G$ (Theorem~\ref{thm:tych}); the results on cones and the answer to Q2 (Section~\ref{sec:cones}); the continuum version of the first countable theorem (Proposition~\ref{prop:fc}); Example~\ref{ex:cZ}; and Theorem~\ref{thm:chain}. We did not find these stated in the sources examined.
\end{enumerate}

\subsection{Sources}\label{sec:sources}
We did not have access to Davison's thesis \cite{davison56}, which Lavallee cites for the notion of curve space and for the local arcwise connectedness of quasi-curve spaces, and we did not consult \cite{davison59} directly. Results of Davison are therefore attributed on the authority of \cite{lavallee64}, and every one that we use is proved in Section~\ref{sec:prelim}. Our account of the example of \cite{fst79} relies in part on the description given by van Douwen \cite[pp.~423--424]{vandouwen81}; we use that example only as a published statement (Section~\ref{sec:fst}).

\subsection{Organisation}
Section~\ref{sec:prelim} contains the definitions and the tools: the basic properties of mosaics, images of compact metric spaces, the sequence characterisations of the three classes, and a criterion for local arcwise connectedness. Section~\ref{sec:arcless} constructs $E(Y)$; the construction of its metric is in Appendix~\ref{app:glue}. Section~\ref{sec:G} constructs $G$ and proves its properties. Section~\ref{sec:published} discusses the published examples. Section~\ref{sec:cones} treats cones and Q2. Section~\ref{sec:fc} contains the first countable results, the example separating the two readings of Q1, and the comparison of the conditions.

\section{Preliminaries}\label{sec:prelim}

\subsection{Conventions}
$\N=\{1,2,3,\dots\}$. Following \cite{lavallee64}, ``metric'' in ``compact metric space'' and ``metric continuum'' means metrizable; no particular metric is part of the structure. A \emph{metric continuum} is a nonempty, compact, connected, metrizable space; a \emph{Peano continuum} (Lavallee: Peano space) is a locally connected metric continuum; a \emph{quasi-Peano space} is a compact, locally connected, metrizable space. A \emph{path} in $X$ is a continuous map $[0,1]\to X$; an \emph{arc} is a space homeomorphic to $[0,1]$. Subsets of a space carry the subspace topology unless something else is said. We do not assume spaces to be Hausdorff.

We use the following classical facts. A Hausdorff space is the image of a path if and only if it is a Peano continuum (Hahn--Mazurkiewicz), and every Peano continuum is arcwise connected and locally arcwise connected \cite[Section~31]{willard}, \cite[Section~6.3]{engelking}. Consequently a compact, locally connected, metrizable space has finitely many components, each open and each a Peano continuum, and a nondegenerate Peano continuum contains an arc.

\subsection{Mosaics}
\begin{definition}[{\cite[p.~1327]{lavallee64}}, after Davison]\label{def:mosaic}
A family $\fM=\{(X_\alpha,\cT_\alpha):\alpha\in A\}$ of topological spaces is a \emph{mosaic on a set $X$} if $X=\bigcup_\alpha X_\alpha$ and, for all $\alpha,\beta$ and every $\cT_\alpha$-closed $M\subseteq X_\alpha$, the set $M\cap X_\beta$ is $\cT_\beta$-closed. The \emph{mosaic topology} $\cT$ on $X$ is defined by: $M\subseteq X$ is closed if and only if $M\cap X_\alpha$ is $\cT_\alpha$-closed for every $\alpha$. We call the $X_\alpha$ the \emph{pieces} of the mosaic.

If all pieces are compact metric spaces, $(X,\cT)$ is a \emph{mosaic space}. A mosaic space obtained from a mosaic all of whose pieces are metric continua (quasi-Peano spaces, Peano continua) is an \emph{$m$-continuum space} (a \emph{quasi-curve space}, a \emph{curve space}).
\end{definition}

\begin{lemma}\label{lem:recognise}
Let $\{(X_\alpha,\cT_\alpha)\}$ be spaces whose underlying sets cover $X$. If for all $\alpha,\beta$ the set $X_\alpha\cap X_\beta$ is closed in $X_\alpha$ and in $X_\beta$ and the topologies $\cT_\alpha$, $\cT_\beta$ induce the same topology on it, then the family is a mosaic on $X$. In particular, any family of closed subspaces of a topological space $Z$ which covers a set $X\subseteq Z$ is a mosaic on $X$, and so is any family of $T_1$ spaces whose pairwise intersections are finite.
\end{lemma}
\begin{proof}
If $M$ is $\cT_\alpha$-closed, then $M\cap X_\beta$ is closed in $X_\alpha\cap X_\beta$ for the common induced topology, and $X_\alpha\cap X_\beta$ is $\cT_\beta$-closed; so $M\cap X_\beta$ is $\cT_\beta$-closed.
\end{proof}

The next proposition collects the basic properties of mosaic spaces. Parts (1), (4) and (5) are attributed to Davison in \cite[Section~1]{lavallee64} (part (4) is Theorem~1.2 there).

\begin{proposition}\label{prop:basic}
Let $\fM=\{(X_\alpha,\cT_\alpha)\}$ be a mosaic of compact metric spaces on $X$ and let $\cT$ be its mosaic topology.
\begin{enumerate}[label=\textup{(\arabic*)}]
\item Each $X_\alpha$ is $\cT$-closed, and $\cT$ induces $\cT_\alpha$ on $X_\alpha$.
\item $(X,\cT)$ is a $T_1$ space.
\item A map from $X$ to a topological space is continuous if and only if its restriction to every piece is continuous. Equivalently, the natural map $\bigsqcup_\alpha X_\alpha\to X$ is a quotient map.
\item If $x_k\to x$ in $X$, there is $\alpha$ such that $x\in X_\alpha$ and $x_k\in X_\alpha$ for infinitely many $k$.
\item Convergent sequences in $X$ have unique limits.
\item $X$ is sequential: a set $W\subseteq X$ is open if and only if every sequence converging to a point of $W$ has all but finitely many terms in $W$.
\end{enumerate}
\end{proposition}
\begin{proof}
(1) Taking $M=X_\alpha$ in Definition~\ref{def:mosaic} shows that $X_\alpha\cap X_\beta$ is $\cT_\beta$-closed for all $\beta$, so $X_\alpha$ is $\cT$-closed. If $M\subseteq X_\alpha$ is $\cT_\alpha$-closed, then $M\cap X_\beta$ is $\cT_\beta$-closed for all $\beta$, so $M$ is $\cT$-closed; conversely, if $M'$ is $\cT$-closed then $M'\cap X_\alpha$ is $\cT_\alpha$-closed. Hence the relatively $\cT$-closed subsets of $X_\alpha$ are exactly the $\cT_\alpha$-closed sets.

(2) A singleton meets each piece in the empty set or a point. (3) is immediate from the definition of $\cT$.

(6) ``Only if'' is trivial. If $W$ is not open, then $(X\setminus W)\cap X_\alpha$ is not $\cT_\alpha$-closed for some $\alpha$; since $X_\alpha$ is metrizable there is a sequence in $(X\setminus W)\cap X_\alpha$ converging in $X_\alpha$, hence in $X$, to a point of $W$.

For (4) and (5) we first prove:

\smallskip\noindent$(\ast)$ \emph{If $(x_k)$ is an injective sequence which converges in $X$ to a point $x$ different from all $x_k$, then there are $\alpha$, a point $z\in X_\alpha$ and a subsequence of $(x_k)$ which lies in $X_\alpha$ and converges to $z$ in $X_\alpha$.}

\smallskip\noindent Indeed, the set $S=\{x_k:k\in\N\}$ has $x$ in its closure and $x\notin S$, so $S$ is not closed and $S\cap X_\alpha$ is not $\cT_\alpha$-closed for some $\alpha$. As $X_\alpha$ is metrizable there are $z\in X_\alpha\setminus S$ and a sequence in $S\cap X_\alpha$ converging to $z$ in $X_\alpha$; since $z\notin S$ this sequence takes infinitely many values, so it has a subsequence that is a subsequence of $(x_k)$.

(5) Note that in a $T_1$ space a sequence converging to $x$ takes each value $c\neq x$ only finitely often, because $X\setminus\{c\}$ is a neighbourhood of $x$. Suppose $x_k\to x$ and $x_k\to y$ with $x\ne y$. Then every value is taken only finitely often, and after passing to a subsequence we may assume that $(x_k)$ is injective and avoids $x$ and $y$. Let $\alpha$, $z$ and the subsequence $(x_{k_j})$ be as in $(\ast)$. One of $x,y$ differs from $z$, say $x\ne z$. The set $C=\{x_{k_j}:j\in\N\}\cup\{z\}$ is compact in $X_\alpha$, hence $\cT_\alpha$-closed, hence $\cT$-closed by (1). But $x\notin C$ and $x_{k_j}\to x$, a contradiction.

(4) If $x_k=x$ for infinitely many $k$, take any $\alpha$ with $x\in X_\alpha$. Otherwise discard the terms equal to $x$; every value is then taken finitely often, so there is an injective subsequence, and $(\ast)$ gives $\alpha$, $z$ and a further subsequence lying in $X_\alpha$ and converging to $z$. By (5), $z=x$.
\end{proof}
By (1), the subsequence in (4) converges to $x$ in the space $X_\alpha$ as well.

\begin{lemma}[Images of compact metric spaces]\label{lem:image}
Let $X$ be a sequential $T_1$ space with unique sequential limits, for instance a mosaic space, let $C$ be a compact metrizable space and $f\colon C\to X$ continuous. Then $f(C)$ is closed in $X$ and is, as a subspace, compact and metrizable. In particular every compact metrizable subspace of $X$ is closed.
\end{lemma}
\begin{proof}
Let $B\subseteq C$ be closed. If $y_k=f(c_k)$ with $c_k\in B$ and $y_k\to y$, pass to a subsequence with $c_k\to c\in B$; then $y_k\to f(c)$, and by uniqueness of limits $y=f(c)\in f(B)$. So $f(B)$ contains the limits of its convergent sequences, and is closed because $X$ is sequential. Thus $f\colon C\to f(C)$ is a closed continuous surjection onto a $T_1$ space, and its fibres are compact.

\emph{$f(C)$ is Hausdorff.} For $y\neq y'$ in $f(C)$ choose disjoint open sets $V\supseteq f^{-1}(y)$, $V'\supseteq f^{-1}(y')$ in the normal space $C$. The sets $f(C)\setminus f(C\setminus V)$ and $f(C)\setminus f(C\setminus V')$ are open in $f(C)$, contain $y$ and $y'$ respectively, and are disjoint, since a common point $z$ would have $\emptyset\ne f^{-1}(z)\subseteq V\cap V'$.

\emph{$f(C)$ is second countable.} Let $\mathcal B$ be a countable base of $C$ closed under finite unions. If $y\in W$ with $W$ open in $f(C)$, the compact set $f^{-1}(y)$ is contained in some $B\in\mathcal B$ with $B\subseteq f^{-1}(W)$, and then $y\in f(C)\setminus f(C\setminus B)\subseteq W$. So the sets $f(C)\setminus f(C\setminus B)$, $B\in\mathcal B$, form a countable base. A compact Hausdorff second countable space is metrizable, by Urysohn's theorem.
\end{proof}

\begin{remark}\label{rem:hausimage}
If $X$ is an arbitrary Hausdorff space, $C$ is compact metrizable and $f\colon C\to X$ is continuous, then $f(C)$ is again compact and metrizable. Indeed $f$ is a closed map, because closed subsets of $C$ are compact and compact subsets of a Hausdorff space are closed; the second and third paragraphs of the preceding proof then apply without change. This is the standard fact that a Hausdorff continuous image of a compact metrizable space is metrizable. Lemma~\ref{lem:image} is the version we need for mosaic spaces, which are not assumed Hausdorff.
\end{remark}

\begin{corollary}\label{cor:image}
Let $X$ be a mosaic space.
\begin{enumerate}[label=\textup{(\arabic*)}]
\item The union of two metric continua in $X$ with a common point is a metric continuum.
\item The image of a path in $X$ is a Peano continuum. Hence every path connected subset of $X$ is arcwise connected, and any two of its points lie in a metric continuum contained in it.
\item Let $x\in X$ and let $K_1,K_2,\dots$ be metric continua in $X$ containing $x$ such that every open set containing $x$ contains $K_n$ for all but finitely many $n$. Then $\bigcup_nK_n$ is a metric continuum.
\end{enumerate}
\end{corollary}
\begin{proof}
(1) The union is the image of the topological sum of the two continua; apply Lemma~\ref{lem:image}. (2) The image is compact metrizable by Lemma~\ref{lem:image}, in particular Hausdorff, so it is a Peano continuum by the Hahn--Mazurkiewicz theorem, and Peano continua are arcwise connected. (3) Let $C$ be the one-point compactification of the topological sum $\bigsqcup_nK_n$, a compact metrizable space. Map each $K_n$ into $X$ by inclusion and the point at infinity to $x$. This map is continuous at the point at infinity: the preimage of an open set $O\ni x$ contains that point together with $K_n$ for all but finitely many $n$, and its complement is a closed subset of finitely many $K_n$, hence compact. The image $\bigcup_nK_n$ is compact metrizable by Lemma~\ref{lem:image}, and connected because every $K_n$ contains $x$.
\end{proof}

\begin{proposition}[Davison; {\cite[Theorem~1.6]{lavallee64}}]\label{prop:char}
A topological space is a mosaic space if and only if it is sequential and its convergent sequences have unique limits. Consequently every Hausdorff sequential space is a mosaic space, and every nonempty open subspace of a mosaic space is a mosaic space.
\end{proposition}
\begin{proof}
Necessity is Proposition~\ref{prop:basic}(5),(6). Conversely let $X$ be sequential with unique sequential limits. Singletons contain the limits of their sequences, so $X$ is $T_1$. Let $(x_k)$ converge to $x$ and put $S=\{x\}\cup\{x_k:k\in\N\}$. We claim that every $S'\subseteq S$ with $x\in S'$ is closed. Let a sequence in $S'$ converge to $y$. If it takes some value infinitely often, $y$ is that value, by uniqueness of limits. Otherwise it has a subsequence which is also a subsequence of $(x_k)$, so $y=x$. In both cases $y\in S'$, and $S'$ is closed because $X$ is sequential. It follows that every point of $S$ other than $x$ is isolated in $S$, while every open set containing $x$ contains all but finitely many points of $S$; so $S$ is a closed subspace which is finite or homeomorphic to a convergent sequence with its limit, in particular compact and metrizable. The family of all such sets $S$ covers $X$ and is a mosaic by Lemma~\ref{lem:recognise}. Its mosaic topology contains the topology of $X$; and if $M$ meets every $S$ in a relatively closed set, then $M$ contains the limits of its convergent sequences, so $M$ is closed in $X$. Thus the two topologies coincide.

A Hausdorff space has unique sequential limits. If $U$ is open in a mosaic space $X$, limits in $U$ are unique, and $U$ is sequential: if every sequence in $U$ converging to a point of $W\subseteq U$ is eventually in $W$, the same holds for sequences in $X$, because they are eventually in $U$; so $W$ is open in $X$.
\end{proof}

The following lemma is the device by which a class of mosaic spaces is recognised from its convergent sequences.

\begin{lemma}\label{lem:generate}
Let $X$ be a mosaic space and $\cK$ a family of compact metrizable subspaces of $X$ which covers $X$. Then $\cK$ is a mosaic on $X$, and its mosaic topology is the topology of $X$ if and only if every convergent sequence in $X$ has a subsequence which lies, together with its limit, in a member of $\cK$.
\end{lemma}
\begin{proof}
The members of $\cK$ are closed (Lemma~\ref{lem:image}), so $\cK$ is a mosaic by Lemma~\ref{lem:recognise}; its mosaic topology $\cT'$ contains the topology $\cT$ of $X$. If $\cT'=\cT$, the condition holds by Proposition~\ref{prop:basic}(4). Conversely assume the condition and let $M$ be $\cT'$-closed. If $y_k\in M$ and $y_k\to y$, some $K\in\cK$ contains $y$ and a subsequence; as $M\cap K$ is closed in $K$, $y\in M$. Since $X$ is sequential, $M$ is $\cT$-closed.
\end{proof}

\begin{corollary}\label{cor:classes}
Let $X$ be a mosaic space.
\begin{enumerate}[label=\textup{(\arabic*)}]
\item \textup{(Lavallee \cite[Theorem~2.1]{lavallee64})} $X$ is an $m$-continuum space if and only if every convergent sequence in $X$ has a subsequence which lies, together with its limit, in a metric continuum contained in $X$.
\item \textup{(cf.\ Lavallee \cite[Theorem~6.1]{lavallee64})} The following are equivalent: \textup{(a)} $X$ is a quasi-curve space; \textup{(b)} $X$ is a curve space; \textup{(c)} every convergent sequence in $X$ has a subsequence which lies, together with its limit, in a Peano continuum contained in $X$; \textup{(d)} every convergent sequence in $X$ has a subsequence which lies, together with its limit, in the image of a path.
\end{enumerate}
\end{corollary}
\begin{proof}
(1) If $X$ is the mosaic space of a mosaic of metric continua, the pieces are subspaces of $X$ and Proposition~\ref{prop:basic}(4) gives the condition. Conversely apply Lemma~\ref{lem:generate} to the family of all metric continua in $X$, which contains the singletons.

(2) (b)$\Rightarrow$(a) is trivial. (a)$\Rightarrow$(c): by Proposition~\ref{prop:basic}(4) a subsequence and the limit $x$ lie in a quasi-Peano piece $Q$. The component of $x$ in $Q$ is open in $Q$ and is a Peano continuum; it contains $x$ and all but finitely many terms of the subsequence. (c)$\Rightarrow$(b): apply Lemma~\ref{lem:generate} to the family of all Peano continua in $X$. (c)$\Leftrightarrow$(d): Peano continua are images of paths, and images of paths in $X$ are Peano continua by Corollary~\ref{cor:image}(2).
\end{proof}

\subsection{Spaces determined by paths}
Franklin and Smith-Thomas \cite{fst79} and van Douwen \cite{vandouwen81} work with a notion defined without reference to mosaics. In van Douwen's formulation \cite[p.~423]{vandouwen81}, a family $\mathcal P$ of subsets of a space $X$ \emph{determines} $X$ if a set $F\subseteq X$ is closed whenever $F\cap P$ is closed in $P$ for every $P\in\mathcal P$, and $X$ is \emph{path-determined} if it is determined by the family of all images of paths in $X$.

\begin{proposition}\label{prop:pathdet}
For a Hausdorff space $X$ the following are equivalent:
\begin{enumerate}[label=\textup{(\alph*)}]
\item $X$ is path-determined;
\item a set $F\subseteq X$ is closed whenever $f^{-1}(F)$ is closed for every path $f$ in $X$;
\item $X$ is a quotient of a topological sum of copies of $[0,1]$;
\item $X$ is a curve space;
\item $X$ is a quasi-curve space.
\end{enumerate}
\end{proposition}
\begin{proof}
Let $f$ be a path in the Hausdorff space $X$ and $P$ its image. Then $f\colon[0,1]\to P$ is a closed map, hence a quotient map, so $f^{-1}(F)$ is closed if and only if $F\cap P$ is closed in $P$. This gives (a)$\Leftrightarrow$(b). If (b) holds, the map from the sum of copies of $[0,1]$ indexed by all paths is a quotient map onto $X$; conversely, if $q\colon\bigsqcup_jI_j\to X$ is a quotient map and $F$ has closed preimage under the paths $q|I_j$, then $F$ is closed. So (b)$\Leftrightarrow$(c).

(a)$\Rightarrow$(d): images of paths in $X$ are compact and metrizable (Remark~\ref{rem:hausimage}). If they determine $X$, then a set containing the limits of its convergent sequences meets each of them in a closed set and is therefore closed; so $X$ is sequential, hence a mosaic space (Proposition~\ref{prop:char}). The images of paths are Peano continua, they form a mosaic by Lemma~\ref{lem:recognise}, and its mosaic topology is the topology of $X$ by (a). (d)$\Rightarrow$(a): every Peano piece is the image of a path and is a subspace, so the images of paths determine $X$. (d)$\Leftrightarrow$(e) is Corollary~\ref{cor:classes}(2).
\end{proof}

\subsection{Local properties}
\begin{definition}[{\cite[Definitions~4.1 and 5.1]{lavallee64}}]\label{def:local}
Two points of a space $Y$ are \emph{$m$-strongly connected in $Y$} if some metric continuum contained in $Y$ contains both, and $Y$ is \emph{$m$-strongly connected} if any two of its points are. A space $X$ is \emph{locally $m$-strongly connected} if every open set containing a point $x$ contains an open $m$-strongly connected set containing $x$.
\end{definition}
In the same way $X$ is \emph{locally connected}, \emph{locally path connected}, \emph{locally arcwise connected} if every open set containing a point $x$ contains an open connected, path connected, arcwise connected set containing $x$; here a set is arcwise connected if any two distinct points of it are the end points of an arc contained in it. For metric spaces, Curtis \cite{curtis80} calls $m$-strongly connected sets \emph{continuum-connected} and locally $m$-strongly connected spaces \emph{locally continuum-connected}. We use the forms with open sets, which are the strongest in common use, so that our counterexamples satisfy the hypotheses in every usual sense.

For a mosaic space, Corollary~\ref{cor:image}(2) gives
\begin{equation}\label{eq:localimpl}
\begin{aligned}
\text{locally path connected}&\iff\text{locally arcwise connected}\\
&\ \Longrightarrow\ \text{locally $m$-strongly connected}\ \Longrightarrow\ \text{locally connected}.
\end{aligned}
\end{equation}

\begin{proposition}[Lavallee {\cite[Theorem~5.2]{lavallee64}}]\label{prop:lav52}
Let $X$ be a mosaic space every nonempty open subspace of which is an $m$-continuum space. Then $X$ is locally $m$-strongly connected.
\end{proposition}
\begin{proof}
Let $O$ be open and $x\in O$, and let $C$ be the set of points of $O$ which are $m$-strongly connected to $x$ in $O$. If $K\subseteq O$ is a metric continuum containing $x$, then $K\subseteq C$. Hence any two points $y,y'\in C$ lie in a set $K\cup K'\subseteq C$ which is a metric continuum by Corollary~\ref{cor:image}(1); so $C$ is $m$-strongly connected. To see that $C$ is open, let $z\in C$ and $z_k\to z$, and suppose $z_k\notin C$ for infinitely many $k$. These $z_k$ are eventually in $O$ and converge to $z$ in the $m$-continuum space $O$, so by Corollary~\ref{cor:classes}(1) infinitely many of them lie with $z$ in a metric continuum $L\subseteq O$. If $K\subseteq O$ is a metric continuum containing $x$ and $z$, then $K\cup L\subseteq C$, a contradiction. By Proposition~\ref{prop:basic}(6), $C$ is open.
\end{proof}

\begin{proposition}\label{prop:qcopen}
Every nonempty open subspace of a quasi-curve space is an $m$-continuum space.
\end{proposition}
\begin{proof}
Let $U$ be open in the quasi-curve space $X$ and let $y_k\to y$ in $U$. By Corollary~\ref{cor:classes}(2) a subsequence lies with $y$ in a Peano continuum $P\subseteq X$. In the locally connected, regular space $P$ there is a connected open set $W\ni y$ whose closure in $P$ lies in $U$. That closure is a metric continuum in $U$ containing $y$ and all but finitely many terms of the subsequence. Since $U$ is a mosaic space (Proposition~\ref{prop:char}), Corollary~\ref{cor:classes}(1) applies.
\end{proof}

The next lemma reduces local arcwise connectedness of a mosaic space to a statement about the pieces that are not locally connected. In the examples below there is essentially one such piece.

\begin{lemma}\label{lem:lac}
Let $X$ be the mosaic space of a mosaic $\fM$ of compact metric spaces. Suppose that for every open $O\subseteq X$, every piece $X_\alpha$ that is not locally connected, every $z\in O\cap X_\alpha$ and every sequence $(z_k)$ in $X_\alpha$ converging to $z$, the point $z_k$ lies in the path component of $z$ in $O$ for all but finitely many $k$. Then the path components of every open subset of $X$ are open and arcwise connected. In particular $X$ is locally arcwise connected.
\end{lemma}
\begin{proof}
Let $C$ be a path component of an open set $O$, let $z\in C$ and $z_k\to z$. Suppose $z_k\notin C$ for infinitely many $k$, and pass to the subsequence of these. By Proposition~\ref{prop:basic}(4) a further subsequence lies, with $z$, in one piece $X_\alpha$ and converges to $z$ there. If $X_\alpha$ is locally connected, it is locally path connected, so $O\cap X_\alpha$ contains a path connected neighbourhood of $z$ in $X_\alpha$, which contains $z_k$ for large $k$ in the subsequence; a path in $X_\alpha$ is a path in $X$, so these $z_k$ are in $C$. If $X_\alpha$ is not locally connected, the hypothesis gives the same conclusion. This contradiction shows, by Proposition~\ref{prop:basic}(6), that $C$ is open. It is arcwise connected by Corollary~\ref{cor:image}(2).
\end{proof}

\begin{corollary}\label{cor:qclac}
Every quasi-curve space is locally arcwise connected, hence locally $m$-strongly connected.
\end{corollary}
\begin{proof}
In a mosaic of quasi-Peano spaces the hypothesis of Lemma~\ref{lem:lac} is vacuous.
\end{proof}
\noindent Lavallee \cite[p.~1333]{lavallee64} attributes the first assertion of Corollary~\ref{cor:qclac} to Davison's thesis.

\section{A metric space without arcs: the answer to Q1}\label{sec:arcless}

\subsection{The idea}
By Proposition~\ref{prop:fc} below, a first countable, locally $m$-strongly connected mosaic space is automatically an $m$-continuum space. So for metric spaces Q1b asks whether small continua joining nearby points can be replaced by small arcs. In a completely metrizable space they can, by the theorem of Mazurkiewicz, Moore and Menger quoted in the introduction. We therefore build an incomplete metric space out of copies of a continuum that contains no arc, using so many copies that any two points at distance $r$ lie in a copy of diameter exactly $r$. Two requirements shape the construction. Every point should be an end point of small copies leading to all nearby points; a single copy cannot be relied on for this, since in the pseudo-arc, for instance, there are arbitrarily close points that lie in no common proper subcontinuum. So copies are attached at \emph{all} pairs of points, and the process is iterated. And we topologise the union by a metric in which each new copy has diameter equal to the distance between its end points; this is precisely what makes open balls $m$-strongly connected.

\subsection{Attaching edges}
Fix a nondegenerate metric continuum $K$ that contains no arc, for example the pseudo-arc or any nondegenerate hereditarily indecomposable continuum \cite{knaster22,moise48,bing48}; such a continuum contains no arc because an arc is decomposable. Fix a compatible metric $\rho$ on $K$ and points $a,b\in K$ with
\[\rho(a,b)=\diam_\rho K=1,\]
which is possible by compactness after rescaling. Let $K^\circ=K\setminus\{a,b\}$; it is nonempty because a nondegenerate continuum is infinite.

\begin{construction}\label{con:plus}
Let $(Y,d)$ be a metric space and $\cE(Y)=\{(x,y)\in Y\times Y:x\ne y\}$. Put
\[Y^+=Y\sqcup\bigl(\cE(Y)\times K^\circ\bigr).\]
For $e=(x,y)\in\cE(Y)$ define $\iota_e\colon K\to Y^+$ by $\iota_e(a)=x$, $\iota_e(b)=y$ and $\iota_e(u)=(e,u)$ for $u\in K^\circ$. Let $E_e=\iota_e(K)$ be the \emph{edge} of $e$, let $E_e^\circ=\iota_e(K^\circ)$, and define a metric $\rho_e$ on $E_e$ by $\rho_e(\iota_eu,\iota_eu')=d(x,y)\,\rho(u,u')$. Thus $\rho_e(x,y)=d(x,y)=\diam_{\rho_e}E_e$, and $E_e\cap Y=\{x,y\}$.
\end{construction}
Edges are indexed by ordered pairs so that no choice of orientation is needed; each pair of points is joined by two edges.

\begin{lemma}\label{lem:plus}
In the situation of Construction~\ref{con:plus} there is a metric $d^+$ on $Y^+$ such that
\begin{enumerate}[label=\textup{(\arabic*)}]
\item $d^+$ coincides with $d$ on $Y$;
\item $d^+$ coincides with $\rho_e$ on $E_e$, for every $e\in\cE(Y)$;
\item if $e=(x,y)$, $u\in E_e^\circ$ and $w\in Y^+\setminus E_e^\circ$, then $d^+(u,w)\ge r(u)$, where $r(u)=\min\{\rho_e(u,x),\rho_e(u,y)\}>0$.
\end{enumerate}
\end{lemma}
The proof is a routine chain argument and is given in Appendix~\ref{app:glue}. The point of the normalisation $\rho(a,b)=1$ is that the edge joining $x$ and $y$ has length exactly $d(x,y)$, so that attaching it creates no shortcut. By (3), $Y$ is closed in $Y^+$ and each $E_e^\circ$ is open in $Y^+$; by (2), each $E_e$ is a subspace homeomorphic to $K$, of diameter $d(x,y)$, containing $x$ and $y$.

\begin{construction}\label{con:EY}
Let $(Y,d_Y)$ be a metric space. Put $(X_0,d_0)=(Y,d_Y)$ and $(X_{n+1},d_{n+1})=(X_n^+,d_n^+)$, and let $E(Y)=\bigcup_nX_n$ with the metric $d=\bigcup_nd_n$, which is well defined by Lemma~\ref{lem:plus}(1).
\end{construction}

\subsection{The theorem}
\begin{theorem}\label{thm:arcless}
Let $(Y,d_Y)$ be a metric space with at least two points which contains no arc. The metric space $(E(Y),d)$ has the following properties.
\begin{enumerate}[label=\textup{(\arabic*)}]
\item $Y$ is a closed subspace of $E(Y)$ and $d$ extends $d_Y$.
\item For all distinct $x,y\in E(Y)$ there is a subspace $E\subseteq E(Y)$ homeomorphic to $K$ with $x,y\in E$ and $\diam E=d(x,y)$. Consequently every open ball of $E(Y)$ is $m$-strongly connected, and $E(Y)$ is connected.
\item $E(Y)$ contains no arc, and every compact, locally connected subspace of $E(Y)$ is finite.
\item $E(Y)$ is a locally $m$-strongly connected mosaic space; it and all of its nonempty open subspaces are $m$-continuum spaces.
\item $E(Y)$ is not a quasi-curve space, and it is not locally arcwise connected.
\end{enumerate}
\end{theorem}
\begin{proof}
\emph{Preliminaries.} Each $X_n$ is closed in $E(Y)$. Indeed, if $u\in E(Y)\setminus X_n$, then $u\in X_{j+1}\setminus X_j$ for a unique $j\ge n$, so $u\in E_e^\circ$ for an edge $e$ attached to $X_j$, and Lemma~\ref{lem:plus}(3), applied to $X_{j+1}=X_j^+$, gives $d(u,w)\ge r(u)>0$ for all $w\in X_j\supseteq X_n$. The same lemma shows that if $e$ is an edge attached to $X_j$, then $E_e^\circ$ is open in $X_{j+1}$. (It is not open in $E(Y)$, and we do not need that.) Part (1) is the case $n=0$.

(2) Choose $n$ with $x,y\in X_n$. The edge $E=E_{(x,y)}\subseteq X_{n+1}$ is homeomorphic to $K$, contains $x$ and $y$, and has diameter $d(x,y)$, by Lemma~\ref{lem:plus}(2) and because $d$ extends $d_{n+1}$. Now let $B=B(c,r)$ be an open ball. For $y\in B\setminus\{c\}$ let $E_y$ be such a set for the pair $c,y$; every $w\in E_y$ satisfies $d(c,w)\le\diam E_y=d(c,y)<r$, so $E_y\subseteq B$. Put $E_c=\{c\}$. For $y,y'\in B$ the set $E_y\cup E_{y'}$ is compact, connected and metrizable, lies in $B$ and contains $y$ and $y'$. Finally $E(Y)$ is connected because any two of its points lie in a connected subset.

(3) Suppose $h\colon[0,1]\to E(Y)$ is an embedding with image $L$. The compact metric space $L$ is the union of the countably many closed sets $L\cap X_n$, so by the Baire category theorem some $L\cap X_n$ has nonempty interior in $L$, and $X_n$ contains an arc. Let $n$ be the least integer such that $X_n$ contains an arc, and let $h'\colon[0,1]\to X_n$ be an embedding with image $L'$. Since $X_0=Y$ contains no arc, $n\ge1$. The closed subset $L'\cap X_{n-1}$ of $L'$ contains no arc, hence has empty interior in $L'$; so there is $t$ with $h'(t)\in X_n\setminus X_{n-1}$, say $h'(t)\in E_e^\circ$ for an edge $e$ attached to $X_{n-1}$. As $E_e^\circ$ is open in $X_n$, there is a nondegenerate closed interval $J\ni t$ with $h'(J)\subseteq E_e^\circ$. Then $h'(J)$ is an arc in $E_e$, which is homeomorphic to $K$: a contradiction.

Let $Q\subseteq E(Y)$ be compact and locally connected. It has finitely many components, each a Peano continuum, and a nondegenerate Peano continuum contains an arc. So every component is a point and $Q$ is finite.

(4) A metric space is Hausdorff and sequential, hence a mosaic space (Proposition~\ref{prop:char}). By (2) the open balls are open $m$-strongly connected sets, so $E(Y)$ is locally $m$-strongly connected. Let $U\subseteq E(Y)$ be open and nonempty and let $y_k\to y$ in $U$. Choose $r>0$ with $B(y,r)\subseteq U$ and discard the finitely many $k$ with $d(y,y_k)\ge r$. For the remaining $k$ let $E_k\subseteq B(y,r)$ be a metric continuum of diameter $d(y,y_k)$ containing $y$ and $y_k$ (with $E_k=\{y\}$ if $y_k=y$). All $E_k$ contain $y$ and $\diam E_k\to0$, so $\bigcup_kE_k$ is compact; it is connected, lies in $U$, and contains $y$ and the remaining $y_k$. By Corollary~\ref{cor:classes}(1), applied to the mosaic space $U$, $U$ is an $m$-continuum space.

(5) Suppose $E(Y)$ is the mosaic space of a mosaic of quasi-Peano spaces. The pieces are compact, locally connected subspaces (Proposition~\ref{prop:basic}(1)), so they are finite by (3). Then every subset of $E(Y)$ meets every piece in a closed set, so every subset is closed and $E(Y)$ is discrete. But $E(Y)$ contains copies of $K$, because $Y$ has two distinct points. This contradiction proves the first assertion. For the second, an arcwise connected set with more than one point contains an arc, so by (3) the only arcwise connected subsets of $E(Y)$ are singletons, which are not open.
\end{proof}

\begin{corollary}[Answer to Q1a and Q1b]\label{cor:Q1}
For $Y=K$, the space $E(K)$ is a metrizable, locally $m$-strongly connected (hence locally connected) $m$-continuum space which is not a quasi-curve space. Moreover every nonempty open subspace of $E(K)$ is an $m$-continuum space.
\end{corollary}
The last sentence says that $E(K)$ satisfies the hypothesis of Lavallee's Theorem~5.2 (Proposition~\ref{prop:lav52}), which is the strongest of the conditions considered in Q1 and Q2.

\begin{remark}\label{rem:complete}
$E(Y)$ is never completely metrizable: a completely metrizable, locally connected space is locally arcwise connected by the Mazurkiewicz--Moore--Menger theorem \cite[Section~6.3]{engelking}, contradicting part (5). The hypothesis that $Y$ contains no arc is necessary for (3), since $Y\subseteq E(Y)$. In the proof of (5) the equality of quasi-curve spaces and curve spaces is not used.
\end{remark}

\begin{remark}[Curtis's example]\label{rem:curtis}
In \cite[Example~1.8]{curtis80} Curtis describes a subset of the plane which is connected and locally continuum-connected but not locally path connected. It is obtained from the closure of the graph of $\sin(1/x)$ by inserting successively smaller copies of that continuum, and the construction is indicated with a figure rather than proved in detail. That example contains arcs, so it does not have property (3). Granting its stated properties, it also answers Q1a and Q1b: a metric space is a mosaic space; a first countable, locally $m$-strongly connected mosaic space is an $m$-continuum space (Proposition~\ref{prop:fc}); and a space that is not locally path connected is not a quasi-curve space (Corollary~\ref{cor:qclac}). We have not found this consequence drawn in the literature, and we have not written out a verification of Curtis's example.
\end{remark}

\section{The space \texorpdfstring{$G$}{G}: the answer to Q3}\label{sec:G}

\subsection{The idea}
A counterexample to Q3 cannot be first countable (Proposition~\ref{prop:fcpath}). We want a point $\infty$ and a sequence $x_n\to\infty$ with two properties.
\begin{enumerate}[label=\textup{(\roman*)},leftmargin=2.4em]
\item \emph{Routes survive.} For every closed set $F$ not containing $\infty$, all but finitely many $x_n$ can be joined to $\infty$ by a path that misses $F$. This is Lemma~\ref{lem:routes}; in a space built as ours is, it is exactly what local arcwise connectedness at $\infty$ requires (Lemma~\ref{lem:lac} and Proposition~\ref{prop:Glac}).
\item \emph{No capture.} No metric continuum contains $\infty$ and infinitely many $x_n$. This is Proposition~\ref{prop:nocapture}.
\end{enumerate}
The two properties pull in opposite directions. The first asks that the routes from the points $x_n$ to $\infty$ be plentiful enough to survive the removal of an arbitrary closed set missing $\infty$; the second forbids any compact metrizable connected set from containing routes for infinitely many $n$. Moreover the space has to be sequential, so its topology near $\infty$ is governed by convergent sequences.

Franklin and Smith-Thomas \cite{fst79} and van Douwen \cite{vandouwen81} resolve this tension with a maximal almost disjoint family, and so do we. Our space is built on the pattern of Arens' space. For each $n$ an arc $R_n$ ends at $x_n$ and carries a sequence of points $(n,i)$, $i\in\N$, converging to $x_n$; we call these points \emph{witnesses}. A maximal almost disjoint family $\cD$ of infinite sets of witnesses, each meeting every $R_n$ in a finite set, is chosen. Each $D\in\cD$ converges to a new point $p_D$, to which its members are joined by a null family of arcs, and the points $p_D$ are joined to $\infty$ by arcs $A_D$ forming a null family at $\infty$. A route from $x_n$ to $\infty$ runs along $R_n$ to a witness, from there to some $p_D$, and along $A_D$. Maximality of $\cD$ is used twice, in opposite directions: a closed set missing $\infty$ cannot meet every route from $x_n$ for infinitely many $n$ (Lemma~\ref{lem:routes}), and a closed connected set containing $\infty$ and infinitely many $x_n$ must contain $p_D$ for uncountably many $D$, whereas the points $p_D$ form a discrete set (Proposition~\ref{prop:nocapture}).

\subsection{Null fans}
For a nonempty countable set $I$ let $\mathrm{Fan}(I)=\{o\}\cup(I\times(0,1])$, and write $(i,0)=o$ for every $i\in I$. Fix an injection $\nu\colon I\to\N$ and define
\[\varphi\bigl((i,t),(j,s)\bigr)=\begin{cases}|t-s|/\nu(i)&\text{if }i=j,\\ t/\nu(i)+s/\nu(j)&\text{if }i\ne j.\end{cases}\]
The map $(i,t)\mapsto(t/\nu(i))\mathbf e_i$ into $\ell^1(I)$ is distance preserving, so $\varphi$ is a metric. The sets $S_i=\{(i,t):0\le t\le1\}$ are the \emph{spines}, $o$ is the \emph{centre}, and $(i,1)$ is the \emph{free end} of $S_i$. The following facts are immediate from the formula; note that the ball of radius $r$ about $o$ contains $S_i$ whenever $1/\nu(i)<r$.
\begin{enumerate}[label=\textup{(F\arabic*)}]
\item A set containing $o$ is a neighbourhood of $o$ if and only if it contains all but finitely many spines and, for each of the remaining spines $S_i$, a set $\{(i,t):t<\eta_i\}$ with $\eta_i>0$. In particular the topology does not depend on $\nu$.
\item $\mathrm{Fan}(I)$ is a Peano continuum. Each spine is an arc, and $S_i\setminus\{o\}$ is open.
\item If $y_k\in S_{i_k}$ and the indices $i_k$ are pairwise distinct, then $y_k\to o$.
\item If $\emptyset\ne J\subseteq I$, then $\mathrm{Fan}(J)$ with its own fan topology is a closed subspace of $\mathrm{Fan}(I)$.
\end{enumerate}

\subsection{Construction}
\begin{construction}\label{con:G}
Let $M=\N\times\N$. For $D\subseteq M$ and $n\in\N$ let $D_n=\{i\in\N:(n,i)\in D\}$ be the \emph{section} of $D$ at $n$. Let
\[\cC=\{D\subseteq M:D\text{ is infinite and }D_n\text{ is finite for every }n\}.\]
Two sets are \emph{almost disjoint} if their intersection is finite. Using Zorn's lemma, fix a family $\cD\subseteq\cC$ of pairwise almost disjoint sets which is maximal with this property. Maximality means:
\begin{equation}\label{eq:max}
\text{for every }Y\in\cC\text{ there is }D\in\cD\text{ such that }D\cap Y\text{ is infinite.}
\end{equation}

The underlying set of $G$ is the disjoint union of a point $\infty$; points $p_D$, one for each $D\in\cD$; the set $\N\times[0,1]$; the set $\cD\times(0,1)$; and the set of all triples $(D,m,t)$ with $D\in\cD$, $m\in D$ and $0<t<1$. We use the following parametrisations.
\begin{itemize}[leftmargin=2em]
\item $R_n(t)=(n,t)$ for $t\in[0,1]$. Put $x_n=R_n(0)$, and identify each $m=(n,i)\in M$ with the point $R_n(1/i)$. These points of $M$ are the \emph{witnesses}; the witnesses of \emph{row $n$} are the points $(n,i)$, $i\in\N$.
\item $a_D\colon[0,1]\to G$, for $D\in\cD$: $a_D(0)=p_D$, $a_D(t)=(D,t)$ for $0<t<1$, $a_D(1)=\infty$.
\item $e_{D,m}\colon[0,1]\to G$, for $D\in\cD$ and $m\in D$: $e_{D,m}(0)=p_D$, $e_{D,m}(t)=(D,m,t)$ for $0<t<1$, $e_{D,m}(1)=m$.
\end{itemize}
The pieces are the following compact metric spaces.
\begin{enumerate}[label=\textup{(P\arabic*)}]
\item $R_n=\{n\}\times[0,1]$ with the topology of $[0,1]$, for $n\in\N$.
\item $\Sigma=\{\infty\}\cup\{x_n:n\in\N\}$, with the topology in which each $x_n$ is isolated and the neighbourhoods of $\infty$ are the cofinite sets containing $\infty$.
\item For $D\in\cD$, the set $\Phi_D=\{p_D\}\cup\bigcup_{m\in D}e_{D,m}([0,1])$ with the topology of $\mathrm{Fan}(D)$ transported by the bijection $(m,t)\mapsto e_{D,m}(t)$. Its centre is $p_D$, its spines are the arcs $E_{D,m}=e_{D,m}([0,1])$, and its free ends are the witnesses $m\in D$.
\item For each nonempty countable $\cS\subseteq\cD$, the set $\Lambda_\cS=\{\infty\}\cup\bigcup_{D\in\cS}a_D([0,1])$ with the topology of $\mathrm{Fan}(\cS)$ transported by $(D,t)\mapsto a_D(1-t)$. Its centre is $\infty$, its spines are the arcs $A_D=a_D([0,1])$, and its free ends are the points $p_D$.
\end{enumerate}
Let $\fM_G$ be the family of all these pieces and let $G$ carry the mosaic topology of $\fM_G$; that $\fM_G$ is a mosaic is Lemma~\ref{lem:Gmosaic}.
\end{construction}
The pieces $\Lambda_\cS$ are taken for all countable $\cS$ because the union of all arcs $A_D$, with the topology making them a null family at $\infty$, is compact but not metrizable when $\cD$ is uncountable; it is the mosaic space of its countable subfans.

\begin{lemma}\label{lem:Gmosaic}
$\fM_G$ is a mosaic of compact metric spaces on $G$. Consequently $G$ is a mosaic space, every piece is a closed subspace of $G$, $x_n\to\infty$, and for every $D\in\cD$ each injective sequence of points of $D$ converges to $p_D$.
\end{lemma}
\begin{proof}
The pieces cover $G$. All pieces are metrizable, so by Lemma~\ref{lem:recognise} pairs of pieces with finite intersection need no further attention. The intersections of distinct pieces are as follows.
\begin{center}
\begin{tabular}{@{}ll@{}}
\toprule
pieces & intersection\\
\midrule
$R_n$, $R_{n'}$ $(n\ne n')$ & $\emptyset$\\
$R_n$, $\Sigma$ & $\{x_n\}$\\
$R_n$, $\Phi_D$ & $\{(n,i):i\in D_n\}$, finite because $D\in\cC$\\
$R_n$, $\Lambda_\cS$ & $\emptyset$\\
$\Sigma$, $\Phi_D$ & $\emptyset$\\
$\Sigma$, $\Lambda_\cS$ & $\{\infty\}$\\
$\Phi_D$, $\Phi_{D'}$ $(D\ne D')$ & $D\cap D'$, finite by almost disjointness\\
$\Phi_D$, $\Lambda_\cS$ & $\{p_D\}$ if $D\in\cS$, otherwise $\emptyset$\\
$\Lambda_\cS$, $\Lambda_{\cS'}$ & $\Lambda_{\cS\cap\cS'}$ if $\cS\cap\cS'\ne\emptyset$, otherwise $\{\infty\}$\\
\bottomrule
\end{tabular}
\end{center}
In the last row the intersection is closed in both pieces, and both induce the fan topology on it, by (F4). So $\fM_G$ is a mosaic. The remaining assertions follow from Proposition~\ref{prop:basic}(1): $x_n\to\infty$ in $\Sigma$, and an injective sequence in $D$ consists of free ends of distinct spines of $\Phi_D$, so it converges to $p_D$ by (F3).
\end{proof}

We fix some terminology. The \emph{vertices} of $G$ are the points $\infty$, $x_n$ $(n\in\N)$, $p_D$ $(D\in\cD)$ and the witnesses. The \emph{elementary intervals} are the arcs $R_n([\frac1{i+1},\frac1i])$ for $n,i\in\N$, the spines $E_{D,m}$, and the arcs $A_D$. Each elementary interval is an arc whose two end points are vertices; every point of $G$ that is not a vertex lies in the interior of exactly one elementary interval; and the only vertices that lie on no elementary interval are the points $x_n$. A function on $G$ is continuous if and only if its restrictions to the pieces are (Proposition~\ref{prop:basic}(3)).

\subsection{Closed sets}
The next lemma records how a closed set can meet the pieces. Everything else in this section is derived from it and from \eqref{eq:max}.
\begin{lemma}\label{lem:closed}
Let $F\subseteq G$ be closed.
\begin{enumerate}[label=\textup{(\arabic*)}]
\item If $\infty\notin F$, then $\{n:x_n\in F\}$ is finite and $\cD_F=\{D\in\cD:A_D\cap F\ne\emptyset\}$ is finite.
\item If $D\in\cD$ and $p_D\notin F$, then $F$ meets only finitely many of the spines $E_{D,m}$, $m\in D$. In particular $D\cap F$ is finite.
\item If $x_n\notin F$, then there is $\varepsilon>0$ with $R_n([0,\varepsilon))\cap F=\emptyset$. In particular only finitely many witnesses of row $n$ lie in $F$.
\end{enumerate}
\end{lemma}
\begin{proof}
(1) The first set is finite because $x_n\to\infty$. If $\cD_F$ were infinite, choose distinct $D^{(k)}\in\cD_F$ and $y_k\in A_{D^{(k)}}\cap F$. In the piece $\Lambda_{\{D^{(k)}:k\in\N\}}$ the points $y_k$ lie on distinct spines, so $y_k\to\infty$ by (F3), and $\infty\in F$.
(2) If $F$ met infinitely many spines of $\Phi_D$, points of $F$ on distinct spines would converge to $p_D$ by (F3).
(3) $R_n\setminus F$ is open in $R_n$ and contains $R_n(0)$.
\end{proof}

\subsection{Complete regularity}
We show that $G$ is a Tychonoff space. Besides its own interest, this gives the Hausdorff property and the fact, needed in Proposition~\ref{prop:nocapture}, that the points $p_D$ form a relatively discrete set. The method is to prescribe the values $0$ and $1$ at the vertices and to extend over the elementary intervals.

\begin{lemma}[Binary extension]\label{lem:binary}
Let $F\subseteq G$ be closed, and let $h$ assign a value in $\{0,1\}$ to every vertex of $G$, with value $0$ at every vertex in $F$, such that:
\begin{enumerate}[label=\textup{(\roman*)}]
\item for each $n$, $h(n,i)=h(x_n)$ for all but finitely many $i$;
\item for each $D\in\cD$, $h(m)=h(p_D)$ for all but finitely many $m\in D$;
\item $h(x_n)=h(\infty)$ for all but finitely many $n$, and $h(p_D)=h(\infty)$ for all but finitely many $D\in\cD$.
\end{enumerate}
Then there is a continuous $f\colon G\to[0,1]$ which equals $h$ at every vertex and vanishes on $F$.
\end{lemma}
\begin{proof}
Put $f=h$ at the vertices. On an elementary interval $I$ with end points $u,w$ define $f$ as follows. If $h(u)=h(w)=0$, let $f=0$ on $I$. If $h(u)=h(w)=1$ and $I\cap F=\emptyset$, let $f=1$ on $I$. In the remaining cases consider the function on the closed set $\{u,w\}\cup(I\cap F)$ which equals $h$ at $u$ and $w$ and $0$ on $I\cap F$. It is well defined because an end point with value $1$ is not in $F$, and continuous because such an end point is then an isolated point of that closed set. By Tietze's theorem it extends to a continuous function $I\to[0,1]$; let $f$ be such an extension. In all cases $f$ agrees with $h$ at the end points of $I$, is continuous on $I$, and vanishes on $I\cap F$. So $f$ is well defined on $G$ and vanishes on $F$. We check continuity on each piece.

\emph{On $R_n$.} On $R_n((0,1])$ the function is continuous, being continuous on each of the intervals $R_n([\frac1{i+1},\frac1i])$, which form a locally finite closed cover of $R_n((0,1])$. At $x_n$: if $h(x_n)=0$, then by (i) all elementary intervals of $R_n$ close enough to $x_n$ have both end points with value $0$, so $f=0$ on a neighbourhood of $x_n$ in $R_n$. If $h(x_n)=1$, then $x_n\notin F$, so by Lemma~\ref{lem:closed}(3) some $R_n([0,\varepsilon))$ misses $F$; with (i), all elementary intervals close enough to $x_n$ have both end points with value $1$ and miss $F$, so $f=1$ on a neighbourhood of $x_n$ in $R_n$.

\emph{On $\Phi_D$.} The function is continuous on each spine, and $E_{D,m}\setminus\{p_D\}$ is open in $\Phi_D$ by (F2), so $f$ is continuous at every point other than $p_D$. If $h(p_D)=0$, then by (ii) all but finitely many spines have both end points with value $0$ and carry $f=0$. If $h(p_D)=1$, then $p_D\notin F$, so $F$ meets only finitely many spines (Lemma~\ref{lem:closed}(2)); with (ii), all but finitely many spines have both end points with value $1$, miss $F$, and carry $f=1$. In both cases $f$ is constant with value $h(p_D)$ on all but finitely many spines and continuous on the others, so it is continuous at $p_D$ by (F1).

\emph{On $\Lambda_\cS$.} The same argument at $\infty$, using (iii) and Lemma~\ref{lem:closed}(1) in place of (ii) and Lemma~\ref{lem:closed}(2).

\emph{On $\Sigma$.} By (iii), $f(x_n)=f(\infty)$ for all but finitely many $n$.
\end{proof}

\begin{theorem}\label{thm:tych}
$G$ is a Tychonoff space: it is $T_1$, and for every closed $F\subseteq G$ and every $q\in G\setminus F$ there is a continuous $f\colon G\to[0,1]$ with $f(q)=1$ and $f=0$ on $F$. In particular $G$ is Hausdorff. Moreover every $p_D$ has an open neighbourhood containing no $p_{D'}$ with $D'\ne D$.
\end{theorem}
\begin{proof}
$G$ is $T_1$ by Proposition~\ref{prop:basic}(2). Let $F$ be closed and $q\notin F$. In each of the first four cases below we specify $h$ and verify the hypotheses of Lemma~\ref{lem:binary}; then $f(q)=h(q)=1$.

\emph{Case $q=\infty$.} By Lemma~\ref{lem:closed}(1) the sets $I_F=\{n:x_n\in F\}$ and $\cD_F$ are finite. Let
\[B=(F\cap M)\cup(I_F\times\N)\cup\bigcup\cD_F,\]
and let $h=1$ on $M\setminus B$, $h=0$ on $B$; $h(x_n)=1$ for $n\notin I_F$ and $h(x_n)=0$ for $n\in I_F$; $h(p_D)=1$ for $D\notin\cD_F$ and $h(p_D)=0$ for $D\in\cD_F$; $h(\infty)=1$. Every vertex in $F$ has value $0$: witnesses in $F$ lie in $B$; if $x_n\in F$ then $n\in I_F$; if $p_D\in F$ then $D\in\cD_F$, as $p_D\in A_D$. Condition (i): if $n\in I_F$, the whole row lies in $B$. If $n\notin I_F$, row $n$ contains only finitely many witnesses in $F$ (Lemma~\ref{lem:closed}(3)), it is disjoint from $I_F\times\N$, and it meets $\bigcup\cD_F$ in a finite set because the members of $\cD_F$ have finite sections; so $h=1$ at all but finitely many witnesses of the row. Condition (ii): if $D\in\cD_F$ then $D\subseteq B$. If $D\notin\cD_F$, then $p_D\notin F$, so $D\cap F$ is finite (Lemma~\ref{lem:closed}(2)); $D\cap(I_F\times\N)$ is finite because $D$ has finite sections; and $D\cap\bigcup\cD_F$ is finite by almost disjointness; so $h=1$ at all but finitely many points of $D$. Condition (iii) holds because $I_F$ and $\cD_F$ are finite.

\emph{Case $q=p_{D_0}$.} Let $h=1$ on $D_0\setminus F$ and $h=0$ on the rest of $M$; $h(p_{D_0})=1$; and $h=0$ at $\infty$, at every $x_n$, and at every $p_D$ with $D\ne D_0$. Vertices in $F$ have value $0$. (i) holds because $D_0$ has finite sections. (ii): $D_0\cap F$ is finite by Lemma~\ref{lem:closed}(2), and $D\cap D_0$ is finite for $D\ne D_0$. (iii) holds with the single exception $D_0$.

\emph{Case $q=x_{n_0}$.} Let $h=1$ at the witnesses of row $n_0$ not in $F$ and $h=0$ on the rest of $M$; $h(x_{n_0})=1$; and $h=0$ at all other vertices. Vertices in $F$ have value $0$. (i): row $n_0$ contains only finitely many witnesses in $F$, by Lemma~\ref{lem:closed}(3). (ii) holds because every $D\in\cD$ has a finite section at $n_0$. (iii) holds with the single exception $n_0$.

\emph{Case $q=m\in M$.} Let $h(m)=1$ and $h=0$ at all other vertices; (i)--(iii) hold with at most one exceptional vertex in each.

\emph{Case $q$ not a vertex.} Then $q$ is an interior point of an elementary interval $I$, and $I$ minus its end points is open in $G$, since its trace on each piece is open there. Choose a continuous $b\colon I\to[0,1]$ with $b(q)=1$ which vanishes on $I\cap F$ and on a neighbourhood of each end point, and let $f=b$ on $I$ and $f=0$ elsewhere. On each piece, $f$ is supported in a compact subset of the interior of one arc of that piece, so it is continuous.

For the last assertion apply the case $q=p_D$ with $F=\emptyset$: the function obtained equals $1$ at $p_D$ and $0$ at every other $p_{D'}$, so $\{f>\frac12\}$ is the required neighbourhood.
\end{proof}

\subsection{Local arcwise connectedness}
\begin{lemma}[Routes]\label{lem:routes}
Let $F\subseteq G$ be closed with $\infty\notin F$. Then for all but finitely many $n$ there is a path in $G\setminus F$ from $x_n$ to $\infty$.
\end{lemma}
\begin{proof}
By Lemma~\ref{lem:closed}(1) the set $N_1=\{n:x_n\in F\}$ and the family $\cD_F$ are finite. For $n\notin N_1$ fix, by Lemma~\ref{lem:closed}(3), $\varepsilon_n\in(0,1)$ with $R_n([0,\varepsilon_n))\cap F=\emptyset$, and let $J_n=\{i:1/i<\varepsilon_n\}$, a cofinite subset of $\N$.

Call $n\notin N_1$ \emph{good} if there are $i\in J_n$ and $D\in\cD\setminus\cD_F$ with $(n,i)\in D$ and $E_{D,(n,i)}\cap F=\emptyset$. For a good $n$, concatenating the path $R_n$ restricted to $[0,1/i]$ (from $x_n$ to the witness $(n,i)$), the path $e_{D,(n,i)}$ reversed (from $(n,i)$ to $p_D$) and the path $a_D$ (from $p_D$ to $\infty$) gives a path from $x_n$ to $\infty$. Each of the three is continuous into a piece, hence into $G$. The first avoids $F$ because $1/i<\varepsilon_n$, the second by the choice of $i$ and $D$, the third because $D\notin\cD_F$.

It remains to show that all but finitely many $n\notin N_1$ are good. Suppose the set $N_2$ of those that are not good is infinite. The set $B=\bigcup\cD_F$ has finite sections, because $\cD_F$ is a finite subfamily of $\cC$. For each $n\in N_2$ choose $i_n\in J_n\setminus B_n$, and let $Y=\{(n,i_n):n\in N_2\}$. Then $Y$ is infinite with sections of at most one element, so $Y\in\cC$, and by \eqref{eq:max} there is $D\in\cD$ with $D\cap Y$ infinite. Since $Y\cap B=\emptyset$, this $D$ is not in $\cD_F$; hence $p_D\notin F$, and by Lemma~\ref{lem:closed}(2) $F$ meets only finitely many spines of $\Phi_D$. So there is $(n,i_n)\in D\cap Y$ with $E_{D,(n,i_n)}\cap F=\emptyset$. This $n$ lies in $N_2$ and is good, a contradiction.
\end{proof}

\begin{proposition}\label{prop:Glac}
The path components of every open subset of $G$ are open and arcwise connected. In particular $G$ is locally arcwise connected.
\end{proposition}
\begin{proof}
We apply Lemma~\ref{lem:lac} to the mosaic $\fM_G$. The pieces $R_n$, $\Phi_D$ and $\Lambda_\cS$ are locally connected, by (F2). The remaining piece is $\Sigma$. Let $O\subseteq G$ be open, $z\in O\cap\Sigma$, and $(z_k)$ a sequence in $\Sigma$ converging to $z$. If $z=x_n$, then $z_k=z$ for all large $k$. If $z=\infty$, apply Lemma~\ref{lem:routes} to $F=G\setminus O$: for all large $k$ either $z_k=\infty$ or $z_k=x_n$ for some $n$ for which $x_n$ is joined to $\infty$ by a path in $O$.
\end{proof}

\subsection{The obstruction for continua}
The next lemma is the connectedness half of the obstruction: a connected closed set cannot reach $x_n$ without containing witnesses arbitrarily close to it.

\begin{lemma}\label{lem:cut}
Let $K\subseteq G$ be closed and connected with $\infty\in K$, and suppose $x_n\in K$. Then $K$ contains infinitely many witnesses of row $n$.
\end{lemma}
\begin{proof}
Suppose $K$ contains only finitely many. Choose $\varepsilon_0\in(0,1)$ such that no witness $(n,i)$ with $1/i\le\varepsilon_0$ lies in $K$. For no $\varepsilon\in(0,\varepsilon_0]$ can $R_n([0,\varepsilon])\subseteq K$ hold, because that arc contains witnesses $(n,i)$ with $1/i\le\varepsilon$. So there is $\varepsilon\in(0,\varepsilon_0]$ with $R_n(\varepsilon)\notin K$. For each $i$ with $1/i<\varepsilon$ and each $D\in\cD$ containing $m=(n,i)$, the set $K\cap E_{D,m}$ is closed in the arc $E_{D,m}$ and does not contain its end point $m$; choose $\delta(D,i)\in(0,1)$ such that $e_{D,m}((1-\delta(D,i),1])\cap K=\emptyset$. Let
\[W=R_n([0,\varepsilon))\ \cup\bigcup_{\substack{i:\ 1/i<\varepsilon\\ D\in\cD:\ (n,i)\in D}}e_{D,(n,i)}\bigl((1-\delta(D,i),1]\bigr).\]
We claim that $W$ is open in $G$, by examining its trace on each piece. On $R_n$ the trace is $R_n([0,\varepsilon))$. On $\Phi_D$ it is the union of the sets $e_{D,(n,i)}((1-\delta(D,i),1])$ over the finitely many $i\in D_n$ with $1/i<\varepsilon$; each of these is open in $\Phi_D$ by (F2). (The sets of this form belonging to another $D'$ meet $\Phi_D$ at most in a witness $(n,i)$ with $i\in D_n$, which is already included.) On $\Sigma$ the trace is $\{x_n\}$. On the pieces $R_{n'}$, $n'\ne n$, and $\Lambda_\cS$ it is empty.

By the choice of the numbers $\delta(D,i)$ and of $\varepsilon$,
\[K\cap W=K\cap R_n([0,\varepsilon))=K\cap R_n([0,\varepsilon]).\]
The first expression shows that this set is open in $K$, the last that it is closed in $K$. It contains $x_n$ and does not contain $\infty$, contradicting the connectedness of $K$.
\end{proof}

\begin{proposition}\label{prop:nocapture}
Let $K\subseteq G$ be closed and connected with $\infty\in K$, and suppose that every discrete subspace of $K$ is countable. Then $x_n\in K$ for only finitely many $n$. In particular no metric continuum in $G$ contains $\infty$ and infinitely many of the points $x_n$.
\end{proposition}
\begin{proof}
A metric continuum in $G$ is closed (Lemma~\ref{lem:image}), connected and second countable, and a discrete subspace of a second countable space is countable; so the second statement follows from the first.

Let $\cD_K=\{D\in\cD:D\cap K\text{ is infinite}\}$. If $D\in\cD_K$, then $p_D\in K$: an injective sequence in $D\cap K$ converges to $p_D$ by Lemma~\ref{lem:Gmosaic}, and $K$ is closed. By Theorem~\ref{thm:tych} the set $\{p_D:D\in\cD_K\}$ is a discrete subspace of $K$. Hence $\cD_K$ is countable, and we can write it as the union of an increasing sequence of finite families $\mathcal F_1\subseteq\mathcal F_2\subseteq\cdots$ (all empty if $\cD_K=\emptyset$).

Suppose now that $x_n\in K$ for infinitely many $n$, say for $n_1<n_2<\cdots$. By Lemma~\ref{lem:cut}, $K$ contains infinitely many witnesses of row $n_k$, whereas $\bigcup\mathcal F_k$ contains only finitely many witnesses of that row, since its members have finite sections. Choose a witness $w_k$ of row $n_k$ with $w_k\in K\setminus\bigcup\mathcal F_k$, and let $Y=\{w_k:k\in\N\}$. Then $Y\in\cC$. If $D\in\cD_K$, say $D\in\mathcal F_j$, then $w_k\notin D$ for all $k\ge j$, so $Y\cap D$ is finite. On the other hand, by \eqref{eq:max} there is $D\in\cD$ with $D\cap Y$ infinite; as $Y\subseteq K$, this $D$ belongs to $\cD_K$. This is a contradiction.
\end{proof}

\begin{remark}\label{rem:spread}
The hypothesis on $K$ in Proposition~\ref{prop:nocapture} is that $K$ has countable spread. It cannot be weakened to separability. Take any $D_0\in\cD$, let $N'$ be the infinite set of $n$ with $(D_0)_n\ne\emptyset$, and let $K$ be the closure of $\bigcup_{n\in N'}R_n\cup\Phi_{D_0}\cup A_{D_0}$. The set whose closure is taken is connected, since every $R_n$ with $n\in N'$ meets $\Phi_{D_0}$, and it is separable; so $K$ is closed, connected and separable, and it contains $\infty$ and $x_n$ for all $n\in N'$. By the proof of the proposition, $K$ contains uncountably many of the points $p_D$.
\end{remark}

\subsection{Summary}
\begin{theorem}\label{thm:G}
The space $G$ of Construction~\ref{con:G} has the following properties.
\begin{enumerate}[label=\textup{(\arabic*)}]
\item $G$ is a mosaic space and a Tychonoff space; in particular it is a Hausdorff sequential space.
\item The path components of every open subset of $G$ are open and arcwise connected; thus $G$ is locally arcwise connected, locally $m$-strongly connected and locally connected.
\item $x_n\to\infty$, and no metric continuum in $G$ contains $\infty$ and infinitely many of the points $x_n$.
\item $G$ is not an $m$-continuum space, and therefore not a quasi-curve space; no path in $G$ passes through infinitely many of the points $x_n$.
\end{enumerate}
\end{theorem}
\begin{proof}
(1) Lemma~\ref{lem:Gmosaic} and Theorem~\ref{thm:tych}. (2) Proposition~\ref{prop:Glac} and \eqref{eq:localimpl}. (3) Lemma~\ref{lem:Gmosaic} and Proposition~\ref{prop:nocapture}. (4) If $G$ were an $m$-continuum space, Corollary~\ref{cor:classes}(1) would give a metric continuum containing $\infty$ and infinitely many $x_n$. A quasi-curve space is an $m$-continuum space (Corollary~\ref{cor:classes}). The image of a path is a metric continuum and is closed (Corollary~\ref{cor:image}(2), Lemma~\ref{lem:image}); if it contained infinitely many $x_n$ it would also contain their limit $\infty$.
\end{proof}

\begin{corollary}[Answer to Q3]\label{cor:Q3}
$G$ is a locally arcwise connected mosaic space which is not a quasi-curve space. In the terminology of \cite{vandouwen81}, $G$ is a completely regular, sequential, locally pathwise connected space which is not path-determined.
\end{corollary}
\begin{proof}
Theorem~\ref{thm:G} and Proposition~\ref{prop:pathdet}.
\end{proof}

\begin{remark}\label{rem:Gprops}
(1) $G$ is not first countable at $\infty$, by Theorem~\ref{thm:G} and Proposition~\ref{prop:fc}. (2) $G$ is not a Fr\'echet--Urysohn space: $\infty$ is in the closure of the set $M$ of witnesses, because each $x_n$ is, but no sequence in $M$ converges to $\infty$, since no piece containing $\infty$ contains a witness (Proposition~\ref{prop:basic}(4)). (3) $G$ is a locally $m$-strongly connected mosaic space which is not an $m$-continuum space; so the variant of Q2 in which $X$ is not assumed to be an $m$-continuum space already fails for $G$. (4) The construction uses the axiom of choice through the family $\cD$, as do the constructions in \cite{fst79} and \cite{vandouwen81}. The finiteness of the sections of the members of $\cD$ is what keeps the sets $D\in\cD$ from accumulating at the points $x_n$; it is used in Lemma~\ref{lem:Gmosaic}, in each case of Theorem~\ref{thm:tych}, and in the two applications of \eqref{eq:max}.
\end{remark}

\section{The published examples}\label{sec:published}
In this section we relate the spaces of Franklin and Smith-Thomas and of van Douwen to Q3, and we prove for van Douwen's space the analogue of Proposition~\ref{prop:nocapture}.

\subsection{The example of Franklin and Smith-Thomas}\label{sec:fst}
Franklin and Smith-Thomas \cite{fst79} study the Hausdorff spaces that are quotients of topological sums of copies of the unit interval, that is, by Proposition~\ref{prop:pathdet}, the Hausdorff path-determined spaces. As reported by van Douwen \cite[p.~424]{vandouwen81}, they show that first countable, locally pathwise connected spaces are path-determined, and their main result is a space, called $\HE$, which is Hausdorff, sequential and locally pathwise connected but not path-determined; $\HE$ is neither Fr\'echet nor regular, and the authors ask \cite[4.2, 4.4]{fst79} whether Fr\'echet or regular examples exist. The space $\HE$ is constructed in \cite[Section~3]{fst79} from Arens' space, cones attached over the neighbourhoods of its top point, and a maximal almost disjoint family; the verification of the Hausdorff property is described there as a lengthy case analysis and is not printed.

By Proposition~\ref{prop:char}, Proposition~\ref{prop:pathdet} and \eqref{eq:localimpl}, a Hausdorff, sequential, locally pathwise connected space which is not path-determined is a locally arcwise connected mosaic space which is not a quasi-curve space. Thus the main result of \cite{fst79} is a negative answer to Davison's question, published fifteen years after Lavallee recorded it. We use the properties of $\HE$ only as published statements; we have not reconstructed the omitted verification, and we have not examined whether some metric continuum in $\HE$ captures the sequence that no path captures.

\subsection{Van Douwen's space}\label{sec:vd}
Van Douwen \cite{vandouwen81} constructs a space, called $\SHE$, which is $\sigma$-compact, cosmic, regular, locally pathwise connected and Fr\'echet, and is not path-determined \cite[p.~424]{vandouwen81}. A regular $T_1$ space is Hausdorff and a Fr\'echet space is sequential, so by the same three results $\SHE$ is a locally arcwise connected mosaic space which is not a quasi-curve space: a second published negative answer to Q3, by a space with much better separation and convergence properties than Q3 requires.

We recall the definition \cite[Sections~4 and 8]{vandouwen81}. The underlying set is
\[\SHE=\{\infty\}\cup\N\cup\bigl(\N\times\N\times(0,1)\bigr),\]
and the idea is that for all $k,n\in\N$ the set $\{\infty\}\cup\{(k,n)\}\times(0,1)\cup\{k\}$ is an arc from $\infty$ to $k$, with $(k,n,x)$ close to $k$ for small $x$. Let $\Phi$ be the set of all functions $f\colon\dom f\to\N$ with $\dom f$ an infinite subset of $\N$. Two members of $\Phi$ are \emph{almost disjoint} if they agree at only finitely many points. Let $M_\Phi\subseteq\Phi$ be a maximal family of pairwise almost disjoint functions (van Douwen writes $M$), and let
\[\lambda\colon M_\Phi\to[\tfrac25,\tfrac35]\ \text{ be an injection such that }\lambda[M_\Phi]\text{ contains no uncountable compact subset.}\]
Such an injection exists \cite[p.~435]{vandouwen81}. The condition on the range of $\lambda$ is part of the definition of $\SHE$: van Douwen uses it to complete the proof of the Fr\'echet property at $\infty$ \cite[p.~435]{vandouwen81}. Our own arguments below (Lemma~\ref{lem:SHEbasic} and Proposition~\ref{prop:SHE}) use only that $\lambda$ takes its values in $[\frac25,\frac35]$. For $k\in\N$, $b\in(0,1)$, $f\in M_\Phi$ and $\varepsilon\in(0,\frac15)$ put
\begin{align*}
A_{k,b}&=\{1,\dots,k\}\cup\bigl(\{1,\dots,k\}\times\N\times(0,b]\bigr),\\
B_{f,\varepsilon}&=\bigl\{(k,f(k),z):k\in\dom f,\ |z-\lambda(f)|\le(\tfrac15-\varepsilon)/k\bigr\}.
\end{align*}
The basic neighbourhoods are:
\begin{itemize}[leftmargin=2em]
\item of $(k,n,x)$: the sets $\{(k,n)\}\times(a,b)$ with $0\le a<x<b\le1$;
\item of $k\in\N$: the sets $U(k,a)=\{k\}\cup\bigl(\{k\}\times\N\times(0,a)\bigr)$ with $a\in(0,1)$;
\item of $\infty$: the sets $V(k,b,F,\varepsilon)=\SHE\setminus\bigl(A_{k,b}\cup\bigcup_{f\in F}B_{f,\varepsilon}\bigr)$ with $k\in\N$, $b\in(0,1)$, $F\subseteq M_\Phi$ finite and $\varepsilon\in(0,\frac15)$.
\end{itemize}
A set is open if it contains a basic neighbourhood of each of its points. The intersection of two basic neighbourhoods of a point contains a basic neighbourhood of that point (for $\infty$, note that $A_{k,b}$ increases with $k$ and $b$ and $B_{f,\varepsilon}$ decreases as $\varepsilon$ increases), so the open sets form a topology. For $f\in M_\Phi$ let
\[L_f=\{(k,f(k),\lambda(f)):k\in\dom f\}.\]

Van Douwen proves that $\SHE$ is locally pathwise connected at $\infty$, is not path-determined, and is Fr\'echet at $\infty$; several routine verifications, among them regularity at $\infty$, are left to the reader \cite[p.~433]{vandouwen81}. The argument for the failure of path-determinedness is that a path meeting $\N$ in an infinite set would contain $\infty$ and, being arcwise connected, would contain for infinitely many $k$ a whole arc $\{(k,n_k)\}\times(0,1)$, hence by maximality infinitely many points of some closed discrete set $L_f$. This uses the arcwise connectedness of paths. Since Q2 concerns continua, we need the corresponding statement for arbitrary compact connected sets, and for it we need only the following elementary properties, which we verify directly from the definition.

\begin{lemma}\label{lem:SHEbasic}
\begin{enumerate}[label=\textup{(\arabic*)}]
\item Every basic neighbourhood is an open set.
\item $\SHE$ is a Hausdorff space.
\item The sequence $(k)_{k\in\N}$ converges to $\infty$.
\item For all $k,n$ the map $x\mapsto(k,n,x)$ is a homeomorphism of $(0,1)$ onto an open subspace of $\SHE$.
\item For $k\in\N$ and $a\in(0,1)$ the set $H(k,a)=\{k\}\cup\bigl(\{k\}\times\N\times(0,a]\bigr)$ is closed.
\item For every $f\in M_\Phi$ the set $L_f$ is closed and discrete.
\end{enumerate}
\end{lemma}
\begin{proof}
(1) This is clear for the neighbourhoods of the points $(k,n,x)$ and for $U(k,a)$. Let $V=V(k,b,F,\varepsilon)$ and $y\in V\setminus\{\infty\}$. If $y=k'\in\N$, then $k'>k$, and $U(k',\frac15)\subseteq V$: a point $(k',n,x)$ with $x<\frac15$ is not in $A_{k,b}$, and it is not in any $B_{f,\varepsilon}$ because $|x-\lambda(f)|>\frac25-\frac15=\frac15>(\frac15-\varepsilon)/k'$. If $y=(k',n,x)$, the set of $z\in(0,1)$ with $(k',n,z)\notin V$ is the union of $(0,b]$ (if $k'\le k$) and the finitely many closed intervals $\{z:|z-\lambda(f)|\le(\frac15-\varepsilon)/k'\}$ for those $f\in F$ with $f(k')=n$; this set is closed in $(0,1)$ and does not contain $x$, so an interval about $x$ gives a basic neighbourhood of $y$ inside $V$.

(2) Two points on the same set $\{(k,n)\}\times(0,1)$ have disjoint interval neighbourhoods, and points on different such sets have disjoint basic neighbourhoods. For $k\ne k'$ in $\N$ the sets $U(k,a)$ and $U(k',a)$ are disjoint, and $U(k,a)$ is disjoint from every basic neighbourhood of a point $(k',n,x)$ with $k'\ne k$. The point $k$ and a point $(k,n,x)$ are separated by $U(k,\frac x2)$ and $\{(k,n)\}\times(\frac x2,1)$. Finally $\infty$ is separated from $k$ by $V(k,\frac12,\emptyset,\frac1{10})$ and $U(k,\frac12)$, and from $(k,n,x)$ by $V(k,b,\emptyset,\frac1{10})$ and $\{(k,n)\}\times(0,b)$ for any $b\in(x,1)$, because $U(k,\frac12)\subseteq A_{k,1/2}$ and $\{(k,n)\}\times(0,b)\subseteq A_{k,b}$.

(3) $V(k,b,F,\varepsilon)$ contains every $k'>k$.

(4) The map is continuous and open onto its image because the sets $\{(k,n)\}\times(a,b)$ are a base at each point of the image, and the image is open.

(5) Let $y\notin H(k,a)$. If $y=(k,n,x)$ with $x>a$, then $\{(k,n)\}\times(a,1)$ misses $H(k,a)$. If $y$ is $k'\ne k$ or a point $(k',n,x)$ with $k'\ne k$, every basic neighbourhood of $y$ misses $H(k,a)$. If $y=\infty$, then $V(k,a,\emptyset,\frac1{10})$ misses $H(k,a)\subseteq A_{k,a}$.

(6) $L_f$ meets each set $\{(k,n)\}\times(0,1)$ in at most one point, so by (4) every point of such a set has a neighbourhood meeting $L_f$ in at most that point. The neighbourhood $U(k,\frac25)$ of $k$ misses $L_f$ because $\lambda(f)\ge\frac25$, and $V(1,\frac12,\{f\},\frac1{10})$ misses $L_f\subseteq B_{f,1/10}$.
\end{proof}

We also need the boundary bumping theorem for Hausdorff continua. It is standard \cite[Section~6.1]{engelking}; we include the short proof in the form we use.

\begin{lemma}[Boundary bumping]\label{lem:bump}
Let $K$ be a compact connected Hausdorff space and $U$ a nonempty open subset of $K$ with $U\ne K$. Then every component of $\cl_KU$ meets $\cl_KU\setminus U$.
\end{lemma}
\begin{proof}
Let $A=\cl_KU$ and let $C$ be a component of $A$ with $C\cap(A\setminus U)=\emptyset$. In the compact Hausdorff space $A$ the component $C$ is the intersection of the sets containing $C$ that are open and closed in $A$ \cite[Section~6.1]{engelking}. These sets are closed under finite intersections and $A\setminus U$ is compact, so one of them, say $E$, misses $A\setminus U$; thus $C\subseteq E\subseteq U$. Now $E$ is closed in $A$, hence in $K$; and $E=W\cap A$ for some open $W\subseteq K$, so $E=E\cap U=W\cap U$ is open in $K$. So $E$ is a nonempty open and closed subset of $K$ with $E\subseteq U\ne K$, contradicting the connectedness of $K$.
\end{proof}

\begin{proposition}\label{prop:SHE}
No compact connected subset of $\SHE$ contains infinitely many points of $\N$. Consequently $\SHE$ is not an $m$-continuum space.
\end{proposition}
\begin{proof}
Let $K\subseteq\SHE$ be compact and connected with $K\cap\N$ infinite. By Lemma~\ref{lem:SHEbasic}(2), $K$ is a compact connected Hausdorff space. Fix $k\in K\cap\N$ and write
\[O=U(k,\tfrac45),\qquad H=H(k,\tfrac45)=O\cup\bigl(\{k\}\times\N\times\{\tfrac45\}\bigr).\]
The set $O$ is open and $H$ is closed (Lemma~\ref{lem:SHEbasic}(1),(5)). Hence $K\cap O$ is an open subset of $K$; it contains $k$, and it is not all of $K$ because $K$ contains points of $\N$ other than $k$, which are not in $O$. All closures and boundaries below are taken in the space $K$. Since $K\cap H$ is closed in $K$ and contains $K\cap O$, we have $\cl_K(K\cap O)\subseteq K\cap H$; and since $K\cap O$ is open in $K$, it lies in the interior of its closure. Therefore
\begin{equation}\label{eq:bump}
\partial_K\,\cl_K(K\cap O)\ \subseteq\ \cl_K(K\cap O)\setminus(K\cap O)\ \subseteq\ K\cap(H\setminus O)\ \subseteq\ \{k\}\times\N\times\{\tfrac45\}.
\end{equation}
By Lemma~\ref{lem:bump}, the component $C$ of $k$ in $\cl_K(K\cap O)$ meets $\cl_K(K\cap O)\setminus(K\cap O)$; by \eqref{eq:bump} it contains a point $(k,n,\frac45)$ for some $n\in\N$.

We claim that $\{(k,n)\}\times(0,\frac45]\subseteq C$. Suppose $(k,n,x_0)\notin C$ with $0<x_0<\frac45$. Since $C\subseteq H$, the set $C$ is the union of
\[C_1=C\cap\bigl(\{(k,n)\}\times(x_0,1)\bigr)\qquad\text{and}\qquad C_2=C\setminus\bigl(\{(k,n)\}\times[x_0,\tfrac45]\bigr).\]
$C_1$ is open in $C$ because $\{(k,n)\}\times(x_0,1)$ is open in $\SHE$. $C_2$ is open in $C$ because $\{(k,n)\}\times[x_0,\frac45]$ is compact by Lemma~\ref{lem:SHEbasic}(4), hence closed in the Hausdorff space $\SHE$. The two sets are disjoint, $(k,n,\frac45)\in C_1$ and $k\in C_2$. This contradicts the connectedness of $C$ and proves the claim.

For each $k\in K\cap\N$ let $g(k)$ be the least $n$ for which $\{(k,n)\}\times(0,\frac45]\subseteq K$. Then $g\in\Phi$, with domain $K\cap\N$. By the maximality of $M_\Phi$ there is $f\in M_\Phi$ which agrees with $g$ at infinitely many points $k$ (if $g\in M_\Phi$ take $f=g$). For each such $k$ the point $(k,f(k),\lambda(f))$ lies in $\{(k,g(k))\}\times(0,\frac45]\subseteq K$, because $\lambda(f)\le\frac35$. So $K$ contains infinitely many points of $L_f$. By Lemma~\ref{lem:SHEbasic}(6) these form an infinite closed discrete subset of the compact space $K$, which is impossible.

For the second statement, suppose $\SHE$ were the mosaic space of a mosaic of metric continua. By Lemma~\ref{lem:SHEbasic}(3) and Proposition~\ref{prop:basic}(4), some piece would contain $\infty$ and infinitely many points of $\N$; a piece is a compact connected subspace.
\end{proof}

Combining Proposition~\ref{prop:SHE} with van Douwen's theorem, $\SHE$ is a regular, Fr\'echet, locally arcwise connected mosaic space which is not an $m$-continuum space. The space $G$ of Section~\ref{sec:G} has the same position with respect to Q3, with complete regularity in place of regularity but without the Fr\'echet property. The mechanisms are closely related: graphs of partial functions $\N\to\N$ are subsets of $\N\times\N$ with sections of at most one element. In $\SHE$ the sets marked by the maximal family are closed and discrete and sit on arcs that join the terms of the sequence directly to $\infty$; in $G$ the marked sets converge to additional points $p_D$, and there are no direct arcs.

\section{Cones: the answer to Q2}\label{sec:cones}

\subsection{The idea}
For Q2 we need an $m$-continuum space $X$, locally $m$-strongly connected, with an open subspace $U$ and a sequence in $U$ which is captured by a metric continuum in $X$ but by none inside $U$. Cones provide this in a uniform way. Let $Y$ be a mosaic space with a mosaic of compact metric pieces $X_\alpha$. The cone over each piece is a metric continuum, and we shall prove that these cones form a mosaic whose mosaic topology is the quotient topology of the cone over $Y$ (Proposition~\ref{prop:cone}). Granting this, every convergent sequence in the cone, including one that converges to the vertex, has a subsequence lying with its limit in the cone over a single piece, by the general property of mosaics (Proposition~\ref{prop:basic}(4)); so the cone is an $m$-continuum space whatever $Y$ is. Note that this is not obtained by projecting to the base: a sequence converging to the vertex need not have a convergent projection.

All these capturing continua pass through the vertex. When the vertex is removed, what remains is the product $Y\times[0,1)$, and a continuum in the product projects onto a continuum in $Y$; so a sequence at height $0$ can be captured in the punctured cone only if its projection is captured in $Y$. Hence the cone over a space such as $G$, in which some sequence is not captured, is a candidate for Q2, provided the cone inherits local path connectedness. Both steps require knowing the topology of the cone precisely, and that is the content of Proposition~\ref{prop:cone}.

\subsection{Products with locally compact spaces}
The following classical result of Whitehead (see \cite[Section~3.3]{engelking}) will be used twice, with $L=[0,1]$; we include its short proof.
\begin{lemma}\label{lem:whitehead}
If $q\colon A\to B$ is a quotient map and $L$ is a locally compact Hausdorff space, then $q\times\mathrm{id}_L\colon A\times L\to B\times L$ is a quotient map.
\end{lemma}
\begin{proof}
Let $W\subseteq B\times L$ have open preimage $W'$ in $A\times L$. For a compact $J\subseteq L$ let $F_J=\{b\in B:\{b\}\times J\not\subseteq W\}$. Then $q^{-1}(F_J)$ is the projection to $A$ of the closed subset $(A\times J)\setminus W'$ of $A\times J$. The projection $A\times J\to A$ is a closed map because $J$ is compact, so $q^{-1}(F_J)$ is closed, and $F_J$ is closed because $q$ is a quotient map. Now let $(b,t)\in W$. For any $a\in q^{-1}(b)$ the set $\{s\in L:(b,s)\in W\}$ equals $\{s:(a,s)\in W'\}$, so it is open in $L$; choose a compact neighbourhood $J$ of $t$ inside it. Then $b\notin F_J$, and $(B\setminus F_J)\times\operatorname{int}J$ is an open neighbourhood of $(b,t)$ contained in $W$.
\end{proof}

\subsection{The cone over a mosaic space}
For a nonempty space $Y$ the \emph{cone} $\cone Y$ is the quotient of $Y\times[0,1]$ obtained by collapsing $Y\times\{1\}$ to a point $v$, the \emph{vertex}; $\kappa\colon Y\times[0,1]\to\cone Y$ denotes the quotient map. As a set, $\cone Y=(Y\times[0,1))\cup\{v\}$. If $B$ is a nonempty compact metrizable space, then $\cone B$ is compact, and it is Hausdorff, being obtained from a compact Hausdorff space by collapsing a closed subset. Hence $\cone B$ is metrizable, as a Hausdorff continuous image of the compact metrizable space $B\times[0,1]$ (Remark~\ref{rem:hausimage}). It is connected, since every point is joined to $v$ by a segment; and by the tube lemma the sets $\kappa(B\times(1-\eta,1])$, $\eta>0$, form a neighbourhood base at $v$.

\begin{proposition}\label{prop:cone}
Let $Y$ be the mosaic space of a mosaic $\fM=\{X_\alpha\}$ of nonempty compact metric spaces, and let $P_\alpha=\kappa(X_\alpha\times[0,1])\subseteq\cone Y$.
\begin{enumerate}[label=\textup{(\arabic*)}]
\item The set $U=\cone Y\setminus\{v\}$ is open, and $\kappa$ restricts to a homeomorphism of $Y\times[0,1)$, with the product topology, onto the subspace $U$.
\item Each $P_\alpha$ is a closed subspace of $\cone Y$, and the map $\cone X_\alpha\to P_\alpha$ induced by $\kappa$ is a homeomorphism. In particular each $P_\alpha$ is a metric continuum.
\item $\{P_\alpha\}$ is a mosaic on $\cone Y$ whose mosaic topology is the topology of $\cone Y$. Hence $\cone Y$ is an $m$-continuum space.
\end{enumerate}
\end{proposition}
\begin{proof}
(1) $Y\times[0,1)$ is open in $Y\times[0,1]$ and is a union of fibres of $\kappa$, and $\kappa$ is injective on it. The restriction of a quotient map to an open union of fibres is a quotient map onto its image, which is open; an injective quotient map is a homeomorphism.

(2) By Proposition~\ref{prop:basic}(3) the natural map $q\colon\bigsqcup_\alpha X_\alpha\to Y$ is a quotient map, so by Lemma~\ref{lem:whitehead} the map $q\times\mathrm{id}\colon\bigsqcup_\alpha(X_\alpha\times[0,1])\to Y\times[0,1]$ is a quotient map, and so is $\kappa\circ(q\times\mathrm{id})$. Thus:
\begin{equation}\label{eq:coneclosed}
\text{$M\subseteq\cone Y$ is closed if and only if $\kappa^{-1}(M)\cap(X_\beta\times[0,1])$ is closed in $X_\beta\times[0,1]$ for every $\beta$.}
\end{equation}
Fix $\alpha$. The restriction of $\kappa$ to $X_\alpha\times[0,1]$ is continuous and constant on $X_\alpha\times\{1\}$, so it induces a continuous bijection $g_\alpha\colon\cone X_\alpha\to P_\alpha$. We show that $g_\alpha$ maps closed sets to closed subsets of $\cone Y$; this proves both assertions. Let $C\subseteq\cone X_\alpha$ be closed and let $C'\subseteq X_\alpha\times[0,1]$ be its preimage, a closed set. For any $\beta$,
\[\kappa^{-1}(g_\alpha(C))\cap(X_\beta\times[0,1])=\bigl(C'\cap((X_\alpha\cap X_\beta)\times[0,1])\bigr)\cup T,\]
where $T=X_\beta\times\{1\}$ if $v\in C$ and $T=\emptyset$ otherwise. (If $v\in C$ then $C'\supseteq X_\alpha\times\{1\}$; if $v\notin C$ then $C'$ is disjoint from $X_\alpha\times\{1\}$.) Since $\fM$ is a mosaic, $X_\alpha\cap X_\beta$ is closed in $X_\beta$ and carries the same topology as a subspace of $X_\alpha$ and of $X_\beta$; hence $C'\cap((X_\alpha\cap X_\beta)\times[0,1])$ is closed in $X_\beta\times[0,1]$. So the displayed set is closed, and $g_\alpha(C)$ is closed by \eqref{eq:coneclosed}.

(3) The $P_\alpha$ are closed subspaces covering $\cone Y$, so they form a mosaic by Lemma~\ref{lem:recognise}. By (2) the map $X_\alpha\times[0,1]\to P_\alpha$ is a quotient map, so $M\cap P_\alpha$ is closed in $P_\alpha$ if and only if $\kappa^{-1}(M)\cap(X_\alpha\times[0,1])$ is closed. By \eqref{eq:coneclosed} the mosaic topology of $\{P_\alpha\}$ is the topology of $\cone Y$. The pieces are metric continua.
\end{proof}

In particular $\cone Y$ is a mosaic space, and its topology does not depend on the mosaic chosen for $Y$. We write $\pi\colon U\to Y$ for the projection $\kappa(y,t)\mapsto y$, which is continuous by part (1).

\begin{theorem}[Cone transfer]\label{thm:conetransfer}
Let $Y$ be a nonempty mosaic space, $\cone Y$ its cone and $U=\cone Y\setminus\{v\}$.
\begin{enumerate}[label=\textup{(\arabic*)}]
\item $\cone Y$ is an $m$-continuum space, and it is contractible.
\item If $Y$ is locally path connected, then the path components of open subsets of $\cone Y$ are open; hence $\cone Y$ is locally arcwise connected and locally $m$-strongly connected.
\item If $Y$ is not an $m$-continuum space, then the open subspace $U$ of $\cone Y$ is not an $m$-continuum space.
\item Every nonempty open subspace of $\cone Y$ is an $m$-continuum space if and only if every nonempty open subspace of $Y$ is an $m$-continuum space.
\item If continuous real-valued functions separate the points of $Y$, they separate the points of $\cone Y$; in particular $\cone Y$ is then Hausdorff.
\end{enumerate}
\end{theorem}
\begin{proof}
(1) The first assertion is Proposition~\ref{prop:cone}(3). For the second, define $H'\colon Y\times[0,1]\times[0,1]\to\cone Y$ by $H'(y,t,r)=\kappa(y,t+(1-t)r)$. It is continuous, and it is constant on the fibres of $\kappa\times\mathrm{id}_{[0,1]}$, since $H'(y,1,r)=v$. By Lemma~\ref{lem:whitehead}, $\kappa\times\mathrm{id}_{[0,1]}$ is a quotient map onto $\cone Y\times[0,1]$ with the product topology. Hence $H'$ induces a continuous map $H\colon\cone Y\times[0,1]\to\cone Y$, and $H(\cdot,0)$ is the identity while $H(\cdot,1)$ is constant. The appeal to Lemma~\ref{lem:whitehead} is what gives joint continuity of $H$; continuity on each set $P_\alpha\times[0,1]$ would not suffice by itself.

(2) $Y\times[0,1]$ is locally path connected. Let $O\subseteq\cone Y$ be open and $P$ a path component of $O$. If $a\in\kappa^{-1}(P)$, the path component of $a$ in the open set $\kappa^{-1}(O)$ is mapped by $\kappa$ into $P$. So $\kappa^{-1}(P)$ is a union of path components of $\kappa^{-1}(O)$, hence open, and $P$ is open because $\kappa$ is a quotient map. The rest follows from \eqref{eq:localimpl}, as $\cone Y$ is a mosaic space.

(3) By Corollary~\ref{cor:classes}(1) there is a sequence $y_k\to y$ in $Y$ no subsequence of which lies, together with $y$, in a metric continuum in $Y$. By Proposition~\ref{prop:cone}(1), $\kappa(y_k,0)\to\kappa(y,0)$ in $U$. If $U$ were an $m$-continuum space, Corollary~\ref{cor:classes}(1), applied to the mosaic space $U$, would give a metric continuum $L\subseteq U$ containing $\kappa(y,0)$ and $\kappa(y_k,0)$ for infinitely many $k$. Then $\pi(L)$ is a metric continuum in $Y$ by Lemma~\ref{lem:image}, and it contains $y$ and infinitely many $y_k$: a contradiction.

(4) Suppose every nonempty open subspace of $\cone Y$ is an $m$-continuum space, and let $V\subseteq Y$ be open and nonempty. By Proposition~\ref{prop:cone}(1), $\kappa(V\times[0,1))$ is open in $\cone Y$. If $y_k\to y$ in $V$, some metric continuum $L\subseteq\kappa(V\times[0,1))$ contains $\kappa(y,0)$ and infinitely many $\kappa(y_k,0)$, and $\pi(L)\subseteq V$ is a metric continuum containing $y$ and infinitely many $y_k$. By Corollary~\ref{cor:classes}(1), $V$ is an $m$-continuum space.

Conversely, suppose every nonempty open subspace of $Y$ is an $m$-continuum space, let $O\subseteq\cone Y$ be open and nonempty, and let $z_j\to z$ in $O$. We look for a metric continuum in $O$ containing $z$ and infinitely many $z_j$.

\emph{Case $z=\kappa(y,t)$ with $t<1$.} By Proposition~\ref{prop:cone}(1) there are an open $V\subseteq Y$ and a compact interval $J\subseteq[0,1)$ which is a neighbourhood of $t$ in $[0,1)$, such that $y\in V$ and $\kappa(V\times J)\subseteq O$. For large $j$ we have $z_j=\kappa(y_j,t_j)$ with $y_j\in V$ and $t_j\in J$, and $y_j\to y$ in $V$. Since $V$ is an $m$-continuum space, a metric continuum $L\subseteq V$ contains $y$ and infinitely many $y_j$. Then $\kappa(L\times J)$ is homeomorphic to $L\times J$, hence a metric continuum; it lies in $O$ and contains $z$ and infinitely many $z_j$.

\emph{Case $z=v$.} If $z_j=v$ for infinitely many $j$, the continuum $\{v\}$ will do. Otherwise, by Proposition~\ref{prop:basic}(4) applied to the mosaic $\{P_\alpha\}$, some $P_\alpha$ contains $v$ and infinitely many $z_j$, and these converge to $v$ in $P_\alpha$. Since $P_\alpha$ is homeomorphic to $\cone X_\alpha$, the open set $O\cap P_\alpha$ contains $\kappa(X_\alpha\times(1-\eta,1])$ for some $\eta>0$. The set $\kappa(X_\alpha\times[1-\frac\eta2,1])$ is compact, metrizable as a subspace of $P_\alpha$, and connected because each of its points is joined to $v$ by a segment inside it; it lies in $O$ and contains all but finitely many of those $z_j$.

In both cases Corollary~\ref{cor:classes}(1), applied to the mosaic space $O$, shows that $O$ is an $m$-continuum space.

(5) The function $\tau\colon\cone Y\to[0,1]$ induced by $(y,t)\mapsto t$ separates points at different heights. If $f\colon Y\to\R$ is continuous with $f(y)\ne f(y')$, the function on $\cone Y$ induced by $(u,s)\mapsto(1-s)f(u)$ separates $\kappa(y,t)$ from $\kappa(y',t)$ for $t<1$.
\end{proof}

\begin{corollary}[Answer to Q2]\label{cor:Q2}
Let $G$ be the space of Section~\ref{sec:G}. The cone $\cone G$ is a contractible, Hausdorff $m$-continuum space which is locally arcwise connected, hence locally $m$-strongly connected, and its open subspace $\cone G\setminus\{v\}$ is not an $m$-continuum space. In particular the converse of Lavallee's Theorem~5.2 is false. Moreover $\cone G$ is not a quasi-curve space.
\end{corollary}
\begin{proof}
$G$ is a locally path connected mosaic space which is not an $m$-continuum space, and continuous functions separate its points (Theorem~\ref{thm:G}). Apply Theorem~\ref{thm:conetransfer}(1), (2), (3), (5). The last assertion follows from Proposition~\ref{prop:qcopen}.
\end{proof}

Thus $\cone G$ is at the same time a counterexample to Q1a, Q1b and Q3: even an $m$-continuum space which is locally arcwise connected need not be a quasi-curve space.

\begin{remark}\label{rem:coneSHE}
Theorem~\ref{thm:conetransfer} applies to any locally path connected mosaic space that is not an $m$-continuum space. By Proposition~\ref{prop:SHE} and van Douwen's theorem, $\SHE$ is such a space, so $\cone\SHE$ is another counterexample to Q2. We do not claim that complete regularity passes from $G$ to $\cone G$, or the Fr\'echet property from $\SHE$ to $\cone\SHE$.
\end{remark}

\section{First countable spaces and the comparison of the conditions}\label{sec:fc}

\subsection{First countable spaces}
In the presence of a countable neighbourhood base at the limit, the diagonal choice described in the introduction is available, and Q2 and Q3 have affirmative answers. The statement for paths is the theorem of \cite{fst79} quoted in Section~\ref{sec:fst}; we give a proof because we need it for mosaic spaces that are not assumed Hausdorff and in the sharper form which captures a whole tail.

\begin{proposition}\label{prop:fc}
Let $X$ be a mosaic space, $U\subseteq X$ open and $x\in U$. Suppose $x$ has a countable neighbourhood base and every open set containing $x$ contains an open $m$-strongly connected set containing $x$. Then for every sequence $(x_n)$ in $U$ converging to $x$ there are $N\in\N$ and a metric continuum $L\subseteq U$ containing $x$ and $x_n$ for all $n\ge N$.

Consequently, in a first countable, locally $m$-strongly connected mosaic space every nonempty open subspace is an $m$-continuum space.
\end{proposition}
\begin{proof}
Let $V_1\supseteq V_2\supseteq\cdots$ be open subsets of $U$ forming a neighbourhood base at $x$, and choose open $m$-strongly connected sets $W_k$ with $x\in W_k\subseteq V_k$. Choose integers $N_1<N_2<\cdots$ such that $x_n\in W_k$ whenever $n\ge N_k$, and for $n\ge N_1$ let $k(n)=\max\{k:N_k\le n\}$. Then $x_n\in W_{k(n)}$ and $k(n)\to\infty$. Let $K_n\subseteq W_{k(n)}$ be a metric continuum containing $x$ and $x_n$. If $O$ is open and $x\in O$, choose $k_0$ with $V_{k_0}\subseteq O$; for $n\ge N_{k_0}$ we have $k(n)\ge k_0$ and $K_n\subseteq V_{k(n)}\subseteq V_{k_0}\subseteq O$. By Corollary~\ref{cor:image}(3), $L=\bigcup_{n\ge N_1}K_n$ is a metric continuum. It lies in $U$ and contains $x$ and all $x_n$ with $n\ge N_1$.

The second statement follows from the first by Proposition~\ref{prop:char} and Corollary~\ref{cor:classes}(1).
\end{proof}
Together with Proposition~\ref{prop:lav52} this shows that for first countable mosaic spaces, in particular for metrizable spaces, ``locally $m$-strongly connected'' is equivalent to ``every nonempty open subspace is an $m$-continuum space''. The sets $W_k$ need not be nested in this proof, because every $K_n$ is joined to the same point $x$. In the next proof consecutive terms of the sequence have to be joined to each other, and nesting becomes essential.

\begin{proposition}\label{prop:fcpath}
Let $X$ be a mosaic space and $x\in X$ a point with a countable neighbourhood base such that every open set containing $x$ contains an open path connected set containing $x$. Then for every sequence $(x_n)$ converging to $x$ there are $N\in\N$ and a path $f\colon[0,1]\to X$ with $f(1)=x$ and $f(1-\frac1j)=x_{N+j-1}$ for all $j\in\N$.

Consequently every first countable, locally arcwise connected mosaic space is a quasi-curve space.
\end{proposition}
\begin{proof}
Let $V_1\supseteq V_2\supseteq\cdots$ be open sets forming a neighbourhood base at $x$. Choose open path connected sets $W_k$ recursively such that
\[x\in W_1\subseteq V_1\qquad\text{and}\qquad x\in W_{k+1}\subseteq W_k\cap V_{k+1}\quad(k\ge1);\]
this is possible because $W_k\cap V_{k+1}$ is an open set containing $x$. The sets $W_k$ are nested, which would not follow from $W_k\subseteq V_k$ alone. Choose $N_1<N_2<\cdots$ with $x_n\in W_k$ for all $n\ge N_k$, and let $k(n)=\max\{k:N_k\le n\}$ for $n\ge N_1$; this function is nondecreasing and tends to infinity. For $n\ge N_1$ we have $x_n\in W_{k(n)}$ and $x_{n+1}\in W_{k(n+1)}\subseteq W_{k(n)}$, so there is a path $\gamma_n$ in $W_{k(n)}$ from $x_n$ to $x_{n+1}$. Let $N=N_1$, let $f$ run along $\gamma_{N+j-1}$ on the interval $[1-\frac1j,1-\frac1{j+1}]$ for each $j\in\N$, and put $f(1)=x$. Then $f(1-\frac1j)=x_{N+j-1}$ and $f$ is continuous on $[0,1)$. If $O$ is open and $x\in O$, choose $k_0$ with $V_{k_0}\subseteq O$; for $n\ge N_{k_0}$ the path $\gamma_n$ lies in $W_{k(n)}\subseteq W_{k_0}\subseteq V_{k_0}\subseteq O$, so $f$ maps a neighbourhood of $1$ into $O$. Hence $f$ is continuous.

For the second statement, the first shows that every convergent sequence has a subsequence lying with its limit in the image of a path; apply Corollary~\ref{cor:classes}(2).
\end{proof}

\subsection{Local connectedness without local \texorpdfstring{$m$}{m}-strong connectedness}
The example of this subsection shows that the hypotheses of Q1a and Q1b differ for $m$-continuum spaces. It depends on the following cited fact: there is a connected, locally connected subset $Z$ of the plane, with more than one point, which contains no nondegenerate continuum. This is Example~1.7 of Curtis \cite{curtis80}, who refers to \cite{hockingyoung}: one takes for $Z$ one of two disjoint sets whose union is the plane and each of which meets every nondegenerate continuum. That $Z$ contains no nondegenerate continuum is then immediate; that $Z$ is connected and locally connected rests on a theorem of plane topology, to the effect that a relatively closed subset of an open disc which separates two points of the disc contains a nondegenerate continuum. We use the existence of $Z$ without proof.

\begin{example}\label{ex:cZ}
Let $Z$ be as above. Then $\cone Z$ is an $m$-continuum space which is locally connected but not locally $m$-strongly connected. In particular $\cone Z$ is not a quasi-curve space.
\end{example}
\begin{proof}
$Z$ is a metric space, hence a mosaic space, and $\cone Z$ is an $m$-continuum space by Proposition~\ref{prop:cone}. Let $U=\cone Z\setminus\{v\}$, which is $Z\times[0,1)$ with the product topology.

\emph{Local connectedness.} Let $O\subseteq\cone Z$ be open and $C$ a component of $O$. If $w\in C$ and $w\ne v$, then $O\cap U$ is an open subset of the locally connected space $Z\times[0,1)$ containing $w$, so it contains a connected open set containing $w$, which lies in $C$. Suppose $v\in C$. Let $W$ be the set of points $(z,t)\in\kappa^{-1}(O)$ with $\{z\}\times[t,1]\subseteq\kappa^{-1}(O)$. Then $W\supseteq Z\times\{1\}$, and $W$ is open in $Z\times[0,1]$: if $(z,t)\in W$, the tube lemma, applied to the compact set $\{z\}\times[t,1]$ in the open set $\kappa^{-1}(O)$, gives a neighbourhood $V'$ of $z$ and an interval $J'=(s,1]$ with $s<t$ such that $V'\times J'\subseteq\kappa^{-1}(O)$, and then $V'\times J'\subseteq W$. As $W$ is a union of fibres of $\kappa$, the set $\kappa(W)$ is open in $\cone Z$. It is connected, being a union of segments $\kappa(\{z\}\times[t,1])$ through $v$, and it lies in $O$; so $v\in\kappa(W)\subseteq C$. Hence $C$ is open.

\emph{Failure of local $m$-strong connectedness.} Let $z_0\in Z$ and suppose $V$ is an open $m$-strongly connected set with $\kappa(z_0,0)\in V\subseteq U$. Since $Z$ is connected with more than one point, $z_0$ is not isolated in $Z$, so $V$ contains a point $\kappa(z,0)$ with $z\ne z_0$. A metric continuum $L\subseteq V$ containing $\kappa(z_0,0)$ and $\kappa(z,0)$ projects under $\pi$ onto a compact connected subset of $Z$ containing two points, that is, onto a nondegenerate continuum in $Z$: a contradiction. The last assertion follows from Corollary~\ref{cor:qclac}.
\end{proof}

\subsection{The four conditions}
\begin{theorem}\label{thm:chain}
For an $m$-continuum space $X$ consider the conditions
\begin{itemize}[leftmargin=4em]
\item[\textup{(QC)}] $X$ is a quasi-curve space;
\item[\textup{(H)}] every nonempty open subspace of $X$ is an $m$-continuum space;
\item[\textup{(LmS)}] $X$ is locally $m$-strongly connected;
\item[\textup{(LC)}] $X$ is locally connected.
\end{itemize}
Then $\textup{(QC)}\Rightarrow\textup{(H)}\Rightarrow\textup{(LmS)}\Rightarrow\textup{(LC)}$, and none of the three implications can be reversed. Moreover local arcwise connectedness is implied by \textup{(QC)} and implies \textup{(LmS)}, and it is independent of \textup{(H)}: neither of the two implies the other.
\end{theorem}
\begin{proof}
The implications are Proposition~\ref{prop:qcopen}, Proposition~\ref{prop:lav52} and \eqref{eq:localimpl}. The space $E(K)$ of Corollary~\ref{cor:Q1} satisfies (H) but not (QC). The space $\cone G$ of Corollary~\ref{cor:Q2} satisfies (LmS) but not (H). The space $\cone Z$ of Example~\ref{ex:cZ} satisfies (LC) but not (LmS). All three are $m$-continuum spaces.

Local arcwise connectedness follows from (QC) by Corollary~\ref{cor:qclac} and implies (LmS) by \eqref{eq:localimpl}. The space $\cone G$ is locally arcwise connected and does not satisfy (H); the space $E(K)$ satisfies (H) and is not locally arcwise connected (Theorem~\ref{thm:arcless}(5)).
\end{proof}
For first countable $m$-continuum spaces the picture is simpler: (H) and (LmS) are equivalent by Proposition~\ref{prop:fc}, and (QC) is equivalent to local arcwise connectedness by Proposition~\ref{prop:fcpath} and Corollary~\ref{cor:qclac}; the metrizable space $E(K)$ shows that these two pairs remain distinct.

\appendix
\section{The metric on a space with attached edges}\label{app:glue}
We prove Lemma~\ref{lem:plus}, in the notation of Construction~\ref{con:plus}. Recall that $E_e\cap Y=\{x,y\}$ for $e=(x,y)$. Since $E_e^\circ\cap E_f=\emptyset$ for $e\ne f$, we have $E_e\cap E_f\subseteq Y$ for $e\neq f$.

Call a pair $(u,w)$ of points of $Y^+$ \emph{elementary} if $u,w\in Y$ or $u,w\in E_e$ for some $e$. Its \emph{cost} $c(u,w)$ is $d(u,w)$ in the first case and $\rho_e(u,w)$ in the second. This is well defined: if $u,w\in Y\cap E_e$, then $u,w\in\{x,y\}$, and $\rho_e$ coincides with $d$ on $\{x,y\}$; if $u,w\in E_e\cap E_f$ with $e\ne f$, then $u,w\in Y$ and both $\rho_e(u,w)$ and $\rho_f(u,w)$ equal $d(u,w)$ by the previous case. A \emph{chain} from $u$ to $w$ is a finite sequence $u=z_0,z_1,\dots,z_k=w$ whose consecutive pairs are elementary; its cost is the sum of the costs of these pairs. Any two points are joined by a chain. Let $d^+(u,w)$ be the infimum of the costs of the chains from $u$ to $w$. Then $d^+$ is symmetric, satisfies the triangle inequality, and vanishes on the diagonal.

\smallskip\noindent\emph{Observation.} If $(z,z')$ is an elementary pair and $z\in E_e^\circ$, then $z'\in E_e$; for $z\notin Y$ and $z\notin E_f$ when $f\neq e$.

\smallskip\noindent\emph{Claim 1. For $u,w\in Y$ every chain from $u$ to $w$ has cost at least $d(u,w)$.} We argue by induction on the number of points of the chain outside $Y$. If there are none, all costs are values of $d$ and the triangle inequality applies. Otherwise let $z_i$ be the first point outside $Y$, say $z_i\in E_e^\circ$ with $e=(x,y)$, and let $j>i$ be the least later index with $z_j\in Y$. Applying the Observation repeatedly, $z_{i-1},z_i,\dots,z_j$ all lie in $E_e$, and $z_{i-1},z_j\in E_e\cap Y=\{x,y\}$. Each pair in this part of the chain contains a point of $E_e^\circ$, so its cost is a value of $\rho_e$, and the part costs at least $\rho_e(z_{i-1},z_j)=d(z_{i-1},z_j)$. Replacing it by the single elementary pair $(z_{i-1},z_j)$ gives a chain of no greater cost with fewer points outside $Y$.

Since the chain $u,w$ has cost $d(u,w)$, Claim~1 gives $d^+=d$ on $Y$, which is (1).

\smallskip\noindent\emph{Claim 2. For $u,w\in E_e$ every chain from $u$ to $w$ has cost at least $\rho_e(u,w)$.} If all points of the chain lie in $E_e$, every cost is a value of $\rho_e$ and the triangle inequality applies. Otherwise let $z_i$ be the first and $z_j$ the last point of the chain outside $E_e$. The pair $(z_{i-1},z_i)$ is elementary with $z_{i-1}\in E_e$ and $z_i\notin E_e$; by the Observation $z_{i-1}\notin E_e^\circ$, so $z_{i-1}\in\{x,y\}$, and similarly $z_{j+1}\in\{x,y\}$. The part of the chain from $z_{i-1}$ to $z_{j+1}$ joins two points of $Y$, so by Claim~1 it costs at least $d(z_{i-1},z_{j+1})=\rho_e(z_{i-1},z_{j+1})$. The parts before $z_{i-1}$ and after $z_{j+1}$ lie in $E_e$ and cost at least $\rho_e(u,z_{i-1})$ and $\rho_e(z_{j+1},w)$. The triangle inequality for $\rho_e$ completes the proof of the claim, and (2) follows.

\smallskip\noindent\emph{Claim 3. If $u\in E_e^\circ$ and $w\notin E_e^\circ$, every chain from $u$ to $w$ has cost at least $r(u)$.} Let $z_i$ be the first point of the chain outside $E_e^\circ$. By the Observation $z_i\in E_e\setminus E_e^\circ=\{x,y\}$, and the part of the chain up to $z_i$ lies in $E_e$, so it costs at least $\rho_e(u,z_i)\ge r(u)$. This is (3); and $r(u)>0$ because $u\notin\{x,y\}$.

\smallskip Finally $d^+$ is a metric. Let $u\ne w$. If both lie in $Y$, or both in one $E_e$, then $d^+(u,w)>0$ by (1) or (2). Otherwise one of them, say $u$, lies in some $E_e^\circ$ while $w\notin E_e$, and $d^+(u,w)\ge r(u)>0$ by (3). \qed

\section*{Use of artificial intelligence}
Artificial intelligence was used extensively in the preparation of this work. Claude Fable 5.2 and ChatGPT Astra assisted in the development of the reported results, proof writing, proofreading, literature searches, revision, and mathematical checking, including attempts to identify errors and gaps. These contributions extended beyond language editing: the constructions and arguments presented here, the identification of the relevant published sources, and the critical review of earlier drafts were all obtained with substantial assistance from these systems. Automated checks are not independent human peer review or machine-checked formal verification, and no such review or verification of this article is claimed.
% NOTE FOR THE AUTHOR: if the target journal requires an author-responsibility
% statement (for example, that the author has verified the content and takes
% responsibility for it), it must be added here by the author. None has been
% written on the author's behalf.

\begin{thebibliography}{99}
\bibitem{bing48} R. H. Bing, \emph{A homogeneous indecomposable plane continuum}, Duke Math. J. \textbf{15} (1948), 729--742.
\bibitem{curtis80} D. W. Curtis, \emph{Hyperspaces of noncompact metric spaces}, Compositio Math. \textbf{40} (1980), 139--152. \url{https://www.numdam.org/item/CM_1980__40_2_139_0/}
\bibitem{davison56} W. F. Davison, \emph{Mosaics}, Ph.D. thesis, University of Virginia, 1956. (Cited in \cite{lavallee64}; not consulted.)
\bibitem{davison59} W. F. Davison, \emph{Mosaics of compact metric spaces}, Trans. Amer. Math. Soc. \textbf{91} (1959), 525--546. \url{https://doi.org/10.1090/S0002-9947-1959-0107853-6}
\bibitem{vandouwen81} E. K. van Douwen, \emph{A locally pathwise connected not path-determined Fr\'echet space, or a method of constructing examples}, Topology Proc. \textbf{6} (1981), 423--436. \url{https://topology.nipissingu.ca/tp/reprints/v06/tp06217.pdf}
\bibitem{engelking} R. Engelking, \emph{General Topology}, revised and completed edition, Sigma Series in Pure Mathematics 6, Heldermann Verlag, Berlin, 1989.
\bibitem{franklin65} S. P. Franklin, \emph{Spaces in which sequences suffice}, Fund. Math. \textbf{57} (1965), 107--115. \url{https://doi.org/10.4064/fm-57-1-107-115}
\bibitem{fst79} S. P. Franklin and B. V. Smith-Thomas, \emph{Topologies determined by paths}, Topology Proc. \textbf{4} (1979), 371--384. \url{https://topology.nipissingu.ca/tp/reprints/v04/tp04204s.pdf}
\bibitem{hockingyoung} J. G. Hocking and G. S. Young, \emph{Topology}, Addison-Wesley, Reading, Mass., 1961.
\bibitem{knaster22} B. Knaster, \emph{Un continu dont tout sous-continu est ind\'ecomposable}, Fund. Math. \textbf{3} (1922), 247--286.
\bibitem{lavallee64} L. D. Lavallee, \emph{Mosaics of metric continua and of quasi-Peano spaces}, Pacific J. Math. \textbf{14} (1964), 1327--1333. \url{https://doi.org/10.2140/pjm.1964.14.1327}
\bibitem{moise48} E. E. Moise, \emph{An indecomposable plane continuum which is homeomorphic to each of its nondegenerate subcontinua}, Trans. Amer. Math. Soc. \textbf{63} (1948), 581--594.
\bibitem{willard} S. Willard, \emph{General Topology}, Addison-Wesley, Reading, Mass., 1970.
\end{thebibliography}
\end{document}
